\documentclass[11pt,reqno]{amsart}

% Encoding and fonts
\usepackage[T1]{fontenc}
\usepackage[utf8]{inputenc}
\usepackage{lmodern}

% Math packages
\usepackage{amsmath}
\usepackage{amssymb}
\usepackage{amsfonts}
\usepackage{amsthm}
\usepackage{mathtools}

% Layout and utilities
\usepackage{geometry}
\usepackage{enumitem}
\usepackage{booktabs}
\usepackage{array}
\usepackage{longtable}
\usepackage[hidelinks]{hyperref}

\geometry{margin=1in}
\numberwithin{equation}{section}
\setcounter{tocdepth}{3}
\setcounter{secnumdepth}{3}

% Theorem-like environments
\theoremstyle{plain}
\newtheorem{theorem}{Theorem}[section]
\newtheorem{proposition}[theorem]{Proposition}
\newtheorem{lemma}[theorem]{Lemma}
\newtheorem{corollary}[theorem]{Corollary}
\newtheorem{conjecture}[theorem]{Conjecture}

\theoremstyle{definition}
\newtheorem{definition}[theorem]{Definition}
\newtheorem{experiment}[theorem]{Experiment}
\newtheorem{heuristic}[theorem]{Heuristic}
\newtheorem{question}[theorem]{Question}

\theoremstyle{remark}
\newtheorem{remark}[theorem]{Remark}
\newtheorem{warning}[theorem]{Warning}

\newcommand{\E}{\mathbb E}
\newcommand{\Prob}{\mathbb P}
\newcommand{\Var}{\operatorname{Var}}
\newcommand{\Cov}{\operatorname{Cov}}
\newcommand{\supp}{\operatorname{supp}}
\newcommand{\Corr}{\operatorname{Corr}}
\newcommand{\Om}{\Omega}
\newcommand{\QF}{\mathcal Q}
\newcommand{\one}{\mathbf 1}
\newcommand{\eps}{\varepsilon}
\newcommand{\cO}{\mathcal O}
\newcommand{\cA}{\mathcal A}

\title[Fixed-Mask Permutations of the M\"obius Path]{A Fixed-Mask Permutation Model for the M\"obius Summatory Path\\
Exact Finite Moments, Quadratic Path Mass, and Computational Placement}
\author{}
\date{Version 3}

\begin{document}
\raggedbottom

\begin{abstract}
For a fixed cutoff \(N\), the nonzero values of the M\"obius function form a
finite sign word on the squarefree support.  We study the uniform orbit
obtained by permuting those signs while preserving the support and the two
sign counts.  This is an exchangeable hypergeometric bridge with the same
endpoint \(M(N)\) as the observed word.  For the normalized logarithmic
quadratic path mass
\begin{equation}
\QF_{N_0,N}(w)=\frac{1}{\log(N/N_0)}
\sum_{N_0<n\leq N}\frac{R_w(n)^2}{n^2},
\qquad R_w(n)=\sum_{k\leq n}w(k),
\label{eq:abstract-quadratic-mass}
\end{equation}
we derive exact finite-class formulae for the mean, the covariance kernel of
squared partial sums, and the variance.  We also give an exact bilinear
decomposition, an opposite-sign transposition calculus, the reflection
degeneracy of completed excursions, and the exact first moment of the number
of returns to zero.

For the identified M\"obius data file through \(N=455{,}052{,}511\), the
reported full-cutoff values are
\begin{equation}
\QF_{1000,N}(\mu)=0.029334,\qquad
\E_{\pi}\QF_{1000,N}(\pi)=0.56628,\qquad
\rho_Q(N)=0.0518.
\label{eq:abstract-full-cutoff}
\end{equation}
The inequality between the first two quantities is a finite, data-dependent
computer-assisted statement with a declared forward-error audit; it is not an
asymptotic theorem.  The approximate five-percent scale is largely explained,
under standard additional hypotheses, by the classical spectral second
moment of the Mertens function.  Thus the class-relative finite comparison is
the contribution of the model, whereas the small absolute quadratic scale is
not claimed as a new phenomenon of analytic number theory.  Reported
excursion, conditional-choice, and blind spectral experiments are included
with their observed outcomes and inferential limitations.  No result in the paper
proves the Riemann Hypothesis or the weak Mertens conjecture.
\end{abstract}

\maketitle
\tableofcontents
\bigskip

\section{Purpose, claims, and novelty boundary}
\label{sec:introduction}

Let \(\mu\) denote the M\"obius function and
\begin{equation}
M(x)=\sum_{n\leq x}\mu(n),\qquad
Q(x)=\sum_{n\leq x}\mu(n)^2.
\label{eq:mertens-and-squarefree-count}
\end{equation}
The model asks a finite conditional question: after retaining the exact
squarefree support and the exact numbers of positive and negative values, how
does the order of the M\"obius signs compare with all other admissible orders?
It does not replace multiplicative number theory by a random-sign hypothesis.
It supplies a null class in which the information removed and the information
retained are explicit.

\subsection{What is established}
\label{subsec:established-summary}

The mathematical core consists of identities valid for every finite
fixed-support sign orbit:
\begin{enumerate}[label=(\roman*)]
\item the exact mean of \eqref{eq:abstract-quadratic-mass};
\item a factorization of
\(\Cov(T_j^2,T_k^2)\) and an \(O(m)\) formula for the exact variance;
\item an exact decomposition into a class-invariant diagonal term and an
order-dependent bilinear term;
\item an exact formula for the change under one opposite-sign transposition;
\item the reflection orbit of completed excursions and an exact formula for
the mean number of zero returns;
\item the diffusive, Brownian-bridge critical scale of the null class.
\end{enumerate}
The numerical component places the observed finite M\"obius word in this
class at fourteen cutoffs.  The strongest data statement is the finite ratio
in \eqref{eq:abstract-full-cutoff}, not a limit as \(N\to\infty\).

\subsection{What is new, and what is classical}
\label{subsec:novelty-boundary}

The finite-population bridge, its concentration inequalities, and its
Brownian-bridge limit belong to standard probability theory
\cite{Hoeffding1963,Serfling1974,Sen1972,BardenetMaillard2015}.  The weak
Mertens conjecture, the conditional spectral second moment, and negative
moments of \(\zeta'\) are classical analytic number theory
\cite{Titchmarsh1986,Ng2004,Gonek1989}.  No priority is claimed for a general
finite-population identity in isolation.

General quadratic forms under permutation and quadratic-assignment methods
already possess an exact-moment literature
\cite{Hoeffding1951,HubertSchultz1976,ZhouEtAl2013}.  The contribution
asserted here is therefore narrower: the particular fixed-mask, fixed-count
orbit for the M\"obius word; the factorization and one-pass variance formula
for its nested-max partial-sum kernel; the local ordering calculus; and the
finite placement of the observed word against that exactly specified class.
No claim is made that exact moments of a generic permuted quadratic form are
new.

\begin{remark}[The meaning of the five-percent observation]
\label{rem:five-percent-status}
The finite value \(0.0518\) is not merely a check that a known conjecture
holds on a finite interval.  It compares two different objects: the observed
M\"obius path and the exact mean of its own fixed-mask, fixed-endpoint orbit.
That particular comparison appears to be specific to the present model.
However, the numerator's scale near \(0.029\) is conditionally explained by
the classical spectral constant in Section \ref{sec:analytic-meaning}.
Consequently, the class-relative normalization is the distinctive part; the
small absolute second moment is not presented as an independent discovery.
\end{remark}

\begin{warning}[No uniqueness claim]
\label{warn:no-uniqueness-from-quadratic-mass}
A single quadratic functional cannot select the M\"obius word uniquely.  It
is constant on an exponentially large excursion-reflection orbit and many
other words have smaller values.  Section \ref{sec:selection-route} gives the
exact arithmetic functional that does select \(\mu\), explains why that
selection is tautological from M\"obius inversion, and states what additional
inequality would be needed to connect selection to RH-scale growth.
\end{warning}

\subsection{Statistical status}
\label{subsec:statistical-status}

The permutation ladder, excursion samples, conditional-choice analyses, and
blind spectral fit are reported computations.  Exact zero-hit bounds use the
one-sided Clopper--Pearson construction \cite{ClopperPearson1934}.  The
conditional-choice \(p\)-values and the blind frequency matches remain
exploratory until their raw outputs and complete selection pipelines are
rerun independently.  Appendix \ref{app:reproducibility} specifies the
required reproducibility package without executing it.

\section{The fixed-mask permutation class}
\label{sec:fixed-mask-class}

\subsection{Support, signs, and class hierarchy}
\label{subsec:support-class}

Fix \(N\geq1\) and write
\begin{equation}
A_N=\{n\leq N:\mu(n)\neq0\}
=\{a_1<a_2<\cdots<a_m\},\qquad m=Q(N).
\label{eq:squarefree-support}
\end{equation}
Let \(p\) and \(q\) be the numbers of \(+1\) and \(-1\) values on \(A_N\).
Then
\begin{equation}
p+q=m,\qquad p-q=S=M(N),\qquad
p=\frac{m+S}{2},\quad q=\frac{m-S}{2}.
\label{eq:sign-counts-endpoint}
\end{equation}

\begin{definition}[Fixed-mask and fixed-count classes]
\label{def:permutation-classes}
Define
\begin{align}
\Om_N^{(0)}
&=\{w\in\{-1,0,1\}^N:\supp(w)=A_N\},
\label{eq:omega-zero}\\
\Om_N^{(1)}
&=\left\{w\in\Om_N^{(0)}:
\#\{n:w(n)=+1\}=p\right\}.
\label{eq:omega-one}\\
\Om_N^{(2)}
&=\left\{w\in\Om_N^{(1)}:
\left|\sum_{k\leq n}w(k)\right|=|M(n)|
\text{ for every }n\leq N\right\}.
\label{eq:omega-two}
\end{align}
Equivalently, \(\Om_N^{(1)}\) is the orbit of
\((\mu(1),\ldots,\mu(N))\) under the symmetric group acting on the values at
the ordered support positions.
\end{definition}

The two defining constraints are the zero mask and the sign multiset.  The
endpoint follows from them; it is not a third constraint.  The first two
cardinalities are
\begin{equation}
|\Om_N^{(0)}|=2^m,\qquad
|\Om_N^{(1)}|=\binom{m}{p}.
\label{eq:class-cardinalities}
\end{equation}
The entire function \(Q(n)\), \(n\leq N\), and the quantities \(p,q,S\) are
class invariants.  The intermediate summatory path is not.

For \(w\in\Om_N^{(1)}\), let
\begin{equation}
\xi_j=w(a_j),\qquad
T_j=\sum_{i=1}^{j}\xi_i,\qquad T_0=0.
\label{eq:support-time-path}
\end{equation}
Then
\begin{equation}
R_w(n)=T_{Q(n)},\qquad R_w(N)=T_m=S.
\label{eq:integer-support-path-link}
\end{equation}
If
\begin{equation}
K_N=\#\{1\leq j\leq m:T_j(\mu)=0\},
\label{eq:number-completed-excursions}
\end{equation}
then Proposition \ref{prop:reflection-orbit} proves
\begin{equation}
\Om_N^{(0)}\supseteq\Om_N^{(1)}\supseteq\Om_N^{(2)},
\qquad
2^m\geq\binom{m}{p}\geq2^{K_N},
\qquad |\Om_N^{(2)}|=2^{K_N}.
\label{eq:class-hierarchy}
\end{equation}
Both inequalities are strict for the reported data at \(N=5\cdot10^6\),
where the base-ten logarithms of the middle and final cardinalities are
\(915017.333\ldots\) and \(1810.996\ldots\), respectively.  At each level,
the uniformization discards the divisor relations except for the invariants
explicitly frozen by that level.
The uniform measure on \(\Om_N^{(1)}\) is sampling without replacement from
a population containing \(p\) values \(+1\) and \(q\) values \(-1\).  Hence
\((T_j)_{0\leq j\leq m}\) is an exchangeable hypergeometric bridge; the
finite-population moment framework is standard \cite{Cochran1977}.

\subsection{Exact exchangeable moments}
\label{subsec:exchangeable-moments}

Assume first that \(m\geq4\).  For distinct indices define
\begin{align}
\mu_1&=\E\xi_a=\frac{S}{m},
\label{eq:first-sign-moment}\\
\mu_2&=\E(\xi_a\xi_b)
=\frac{S^2-m}{m(m-1)},
\label{eq:second-sign-moment}\\
\mu_4&=\E(\xi_a\xi_b\xi_c\xi_d)
=\frac{S^4-6mS^2+8S^2+3m^2-6m}
{m(m-1)(m-2)(m-3)}.
\label{eq:fourth-sign-moment}
\end{align}
The identities follow by expanding powers of
\(\sum_{j=1}^m\xi_j=S\), using \(\xi_j^2=1\), and exploiting
exchangeability.  In particular,
\begin{equation}
\E T_j^2=j+j(j-1)\mu_2.
\label{eq:partial-square-mean}
\end{equation}
For \(m<2\), the second mixed moment is unnecessary and the class is
enumerated directly.  For \(2\leq m<4\), formula
\eqref{eq:second-sign-moment} remains available but fourth-order quantities
are evaluated by direct finite enumeration; no undefined quotient is used.

\begin{theorem}[Factored covariance of squared partial sums]
\label{thm:factored-covariance}
For \(m\geq4\) and \(1\leq j\leq k\leq m\),
\begin{equation}
\Cov(T_j^2,T_k^2)=\phi(j)\psi(k),
\qquad \phi(j)=j(j-1),
\label{eq:factored-covariance}
\end{equation}
where
\begin{equation}
\psi(k)=2+4\mu_2(k-2)+\mu_4(k-2)(k-3)
-\mu_2^2k(k-1).
\label{eq:psi-kernel}
\end{equation}
In particular, \(\psi(m)=0\), as required by the deterministic endpoint
\(T_m=S\).
\end{theorem}

\begin{proof}
Write \(T_k=T_j+U_{j,k}\), expand \(T_j^2T_k^2\), and group monomials by the
number of distinct indices.  Exchangeability replaces all products with two
distinct indices by \(\mu_2\) and all products with four distinct indices by
\(\mu_4\).  Subtracting
\(\E T_j^2\E T_k^2\) and collecting the coefficients gives
\eqref{eq:factored-covariance}--\eqref{eq:psi-kernel}.  The complete
enumeration audit specified in Appendix \ref{app:reproducibility} checks all
pairs for every sign class with \(m\leq8\).
\end{proof}

\subsection{The critical square-root scale}
\label{subsec:critical-scale}

Put
\begin{equation}
U_j=T_j-\frac{jS}{m},\qquad
\sigma_m^2=1-\frac{S^2}{m^2}.
\label{eq:centered-bridge}
\end{equation}
For each fixed \(j\), Hoeffding's comparison theorem for sampling without
replacement gives
\begin{equation}
\Prob\left(\left|U_j\right|\geq\lambda\right)
\leq2\exp\left(-\frac{\lambda^2}{2j}\right).
\label{eq:pointwise-hoeffding}
\end{equation}
Applying a union bound, rather than claiming a sharp maximal inequality,
yields
\begin{equation}
\Prob\left(\max_{1\leq j\leq m}|U_j|>\lambda\right)
\leq2m\exp\left(-\frac{\lambda^2}{2m}\right).
\label{eq:union-maximal-bound}
\end{equation}
Thus, for every fixed \(\eps>0\),
\begin{equation}
\Prob\left(\max_{j\leq m}|U_j|>m^{1/2+\eps}\right)
\leq2m\exp\left(-\frac{m^{2\eps}}{2}\right).
\label{eq:superpolynomial-critical-bound}
\end{equation}

\begin{lemma}[Uniform square integrability of the bridge maximum]
\label{lem:bridge-maximum-ui}
For every nondegenerate fixed-count sign class with \(m\geq3\),
\begin{equation}
\E_{\pi}\left[
\left\{\frac{\max_{0\leq j\leq m}|U_j|}
{\sigma_m\sqrt m}\right\}^{\!2}\right]
\leq32.
\label{eq:bridge-maximum-l2-bound}
\end{equation}
Consequently the normalized maxima are uniformly integrable along every
sequence of such classes.
\end{lemma}

\begin{proof}
Let \(X_i=\xi_i-S/m\), let \(\mathcal F_j\) be the first \(j\) revealed
signs, and put \(h=\lfloor m/2\rfloor\).  Since the unrevealed centered
population has sum \(-U_j\),
\begin{equation}
\E(X_{j+1}\mid\mathcal F_j)=-\frac{U_j}{m-j}.
\label{eq:bridge-conditional-drift}
\end{equation}
It follows that
\begin{equation}
Y_j=\frac{U_j}{m-j},\qquad 0\leq j<m,
\label{eq:bridge-doob-martingale}
\end{equation}
is a martingale.  Doob's \(L^2\) inequality and
\begin{equation}
\E U_h^2=\sigma_m^2\frac{h(m-h)}{m-1}
\label{eq:finite-bridge-variance}
\end{equation}
give
\begin{align}
\E\max_{0\leq j\leq h}U_j^2
&\leq m^2\E\max_{0\leq j\leq h}Y_j^2
\leq4m^2\E Y_h^2
\leq8\sigma_m^2m.
\label{eq:first-half-maximal-bound}
\end{align}
Apply the same argument to the reversed centered sequence.  Since its partial
sums satisfy \(\widetilde U_r=-U_{m-r}\), the second half is bounded by
\begin{equation}
\E\max_{h\leq j\leq m}U_j^2\leq16\sigma_m^2m.
\label{eq:second-half-maximal-bound}
\end{equation}
The square of the full maximum is at most the sum of the two half-maxima.
This gives the stated, slightly relaxed constant \(32\); the finitely many
small boundary cases are checked directly.
\end{proof}

\begin{proposition}[Finite-population bridge limit]
\label{prop:finite-population-fclt}
Suppose \(m\to\infty\) and \(|S|/m\to s_0<1\).  The polygonal interpolation
of
\begin{equation}
\mathbb T_m(t)=\frac{T_{\lfloor mt\rfloor}
-(\lfloor mt\rfloor/m)S}{\sigma_m\sqrt m},
\qquad 0\leq t\leq1,
\label{eq:bridge-process-normalization}
\end{equation}
converges weakly to a standard Brownian bridge \(B^{\circ}\).  Consequently,
\begin{equation}
\frac{1}{\sigma_m\sqrt m}\max_{j\leq m}|U_j|
\Longrightarrow\sup_{0\leq t\leq1}|B^{\circ}(t)|.
\label{eq:maximum-bridge-limit}
\end{equation}
\begin{equation}
\E\max_{j\leq m}|U_j|
\sim\sigma_m\sqrt{\frac{\pi}{2}}\log2\,\sqrt m.
\label{eq:expected-bridge-maximum}
\end{equation}
\end{proposition}

\begin{proof}
The weak convergence is the finite-population functional central limit
theorem; see \cite{Sen1972}.  The continuous-mapping theorem gives
\eqref{eq:maximum-bridge-limit}.  Lemma \ref{lem:bridge-maximum-ui} supplies
uniform integrability, so convergence of the first moments follows.  The
limit distribution and its first moment use
\begin{align}
\Prob\left(\sup_t|B^{\circ}(t)|\leq x\right)
&=\sum_{r\in\mathbb Z}(-1)^r e^{-2r^2x^2},
\label{eq:kolmogorov-distribution}\\
\E\sup_t|B^{\circ}(t)|
&=\sqrt{\frac{\pi}{2}}\log2
=0.868731\ldots,
\label{eq:kolmogorov-mean}
\end{align}
as recorded, for example, in \cite{Feller1968}.
\end{proof}

At \(j=\lfloor m/2\rfloor\), the hypergeometric local central limit theorem
gives, for every fixed \(c>0\),
\begin{equation}
\Prob\left(|U_{\lfloor m/2\rfloor}|>c\sigma_m\sqrt m\right)
\longrightarrow2\overline\Phi(2c)>0.
\label{eq:critical-lower-bound}
\end{equation}
Equations \eqref{eq:superpolynomial-critical-bound} and
\eqref{eq:critical-lower-bound} show that \(m^{1/2}\) is the critical finite
scale of the centered class.  This is the same exponent that occurs in the
classical RH criterion, but the logical statements are different.

\begin{warning}[Endpoint and coupling]
\label{warn:endpoint-and-coupling}
The deterministic endpoint satisfies
\begin{equation}
\left|\max_j|T_j|-\max_j|U_j|\right|\leq|S|,
\qquad \max_j|T_j|\geq|S|.
\label{eq:endpoint-obstruction}
\end{equation}
A concentration statement for \(U_j\) does not control \(S=M(N)\).
Moreover, the finite classes at different values of \(N\) are not a single
canonically coupled random sequence.  The law of the iterated logarithm is
therefore not a prediction for one draw from \(\Om_N^{(1)}\); it concerns an
almost-sure envelope along a coupled infinite sequence
\cite{HartmanWintner1941}.
\end{warning}

\begin{warning}[Why typicality does not imply RH]
\label{warn:typicality-not-rh}
At \(N=5\cdot10^6\), a designated member has uniform mass
\begin{equation}
|\Om_N^{(1)}|^{-1}=10^{-915017.333\ldots}.
\label{eq:single-word-class-mass}
\end{equation}
An event containing a \(1-o(1)\) fraction of a finite class need not contain
that designated member.  In addition, the class deliberately discards the
divisor relations that define \(\mu\).  No concentration inequality over the
class transfers by itself to the M\"obius word.
\end{warning}

\section{Normalized logarithmic quadratic path mass}
\label{sec:quadratic-mass}

\subsection{Definition and plateau representation}
\label{subsec:quadratic-mass-definition}

Fix \(1\leq N_0<N\) and put
\begin{equation}
L_N=\log(N/N_0).
\label{eq:logarithmic-span}
\end{equation}

\begin{definition}[Normalized logarithmic quadratic path mass]
\label{def:quadratic-path-mass}
For a finite word \(w\) and its partial sums \(R_w(n)\), define
\begin{equation}
\QF_{N_0,N}(w)
=\frac{1}{L_N}\sum_{N_0<n\leq N}\frac{R_w(n)^2}{n^2}.
\label{eq:quadratic-path-mass}
\end{equation}
This is a positive logarithmically weighted second moment of the summatory
path.  No physical interpretation is assumed.
\end{definition}

For the fixed mask, define the plateau weights
\begin{equation}
w_j=w_j^{(N_0,N)}
=\sum_{\substack{N_0<n\leq N\\Q(n)=j}}\frac{1}{n^2},
\qquad 1\leq j\leq m.
\label{eq:plateau-weights}
\end{equation}
Then the path functional is the quadratic form
\begin{equation}
L_N\QF_{N_0,N}(w)=\sum_{j=1}^{m}w_jT_j^2.
\label{eq:plateau-quadratic-form}
\end{equation}
The weights depend only on the squarefree mask; all order dependence is in
the bridge coordinates \(T_j\).

\subsection{Relation to quadratic forms under permutation}
\label{subsec:permutation-quadratic-forms}

Let \(\xi=(\xi_1,\ldots,\xi_m)^{\mathsf T}\).  Expanding the partial sums in
\eqref{eq:plateau-quadratic-form} gives
\begin{align}
L_N\QF_{N_0,N}(w)
&=\xi^{\mathsf T}\mathbf K_N\xi,
\label{eq:permutation-quadratic-form}\\
(\mathbf K_N)_{rs}
&=\sum_{j=\max(r,s)}^m w_j,
\qquad 1\leq r,s\leq m.
\label{eq:nested-max-kernel}
\end{align}
Thus the model is a quadratic form evaluated on a uniform permutation of a
fixed multiset of signs.  Permutation quadratic forms, quadratic-assignment
statistics, and their exact moments have an established literature
\cite{Hoeffding1951,HubertSchultz1976,ZhouEtAl2013}.  In particular, exact
permutation moments of general quadratic forms are not claimed here as a new
category of result.

The specific feature used below is the nested-max kernel
\eqref{eq:nested-max-kernel}.  For this kernel and the fixed-count sign
population, the covariance of squared prefix sums factorizes on the upper
triangle, reducing the variance to one prefix pass; the same prefix geometry
also yields the local transposition formula in Section
\ref{subsec:transposition-calculus}.  These special simplifications, rather
than the generic quadratic-form representation, delimit the finite
probabilistic contribution.

\subsection{Exact finite-class mean and variance}
\label{subsec:exact-mean-variance}

Define two mask sums
\begin{align}
\Sigma_1(N_0,N)
&=\sum_{N_0<n\leq N}\frac{Q(n)}{n^2}
=\sum_{j=1}^m j w_j,
\label{eq:sigma-one}\\
\Sigma_2(N_0,N)
&=\sum_{N_0<n\leq N}\frac{Q(n)^2}{n^2}
=\sum_{j=1}^m j^2w_j.
\label{eq:sigma-two}
\end{align}
For \(m\geq2\), let
\begin{equation}
A=\frac{m^2-S^2}{m(m-1)}=1-\mu_2,\qquad
B=\frac{S^2-m}{m(m-1)}=\mu_2.
\label{eq:quadratic-reduction-coefficients}
\end{equation}

\begin{theorem}[Exact mean over the complete class]
\label{thm:exact-class-mean}
For the uniform measure on \(\Om_N^{(1)}\),
\begin{align}
\overline{\QF}_N
:=\E_{\pi}\QF_{N_0,N}(\pi)
&=\frac{A\Sigma_1+B\Sigma_2}{L_N}
\label{eq:exact-class-mean}\\
&=\frac{1}{L_N}\sum_{j=1}^m
w_j\{j+\mu_2j(j-1)\}.
\label{eq:exact-class-mean-alternative}
\end{align}
Thus all cutoffs can be evaluated in one streaming pass, because
\(\Sigma_1\) and \(\Sigma_2\) are cumulative while \(A\) and \(B\) depend
only on the terminal values \(m\) and \(S\).
\end{theorem}

\begin{proof}
Substitute \eqref{eq:partial-square-mean} into
\eqref{eq:plateau-quadratic-form}.  The identity
\(\E T_j^2=Aj+Bj^2\) then gives \eqref{eq:exact-class-mean}; expanding
\(A=1-\mu_2\) gives \eqref{eq:exact-class-mean-alternative}.
\end{proof}

\begin{theorem}[Exact variance in one prefix pass]
\label{thm:exact-class-variance}
Assume \(m\geq4\), and define
\begin{equation}
P_k=\sum_{j<k}w_j\phi(j).
\label{eq:variance-prefix}
\end{equation}
Then
\begin{equation}
\Var_{\pi}\!\left(\QF_{N_0,N}(\pi)\right)
=\frac{1}{L_N^2}\sum_{k=1}^{m}
w_k\psi(k)\{w_k\phi(k)+2P_k\}.
\label{eq:exact-class-variance}
\end{equation}
The formula uses \(O(m)\) operations and \(O(1)\) running state once the
weights are streamed.
\end{theorem}

\begin{proof}
Expand the variance of \(\sum_jw_jT_j^2\), use symmetry of covariance, and
apply Theorem \ref{thm:factored-covariance}:
\begin{align}
L_N^2\Var(\QF)
&=\sum_{k=1}^m w_k^2\phi(k)\psi(k)
+2\sum_{1\leq j<k\leq m}w_jw_k\phi(j)\psi(k)
\label{eq:variance-triangular-sum}\\
&=\sum_{k=1}^{m}w_k\psi(k)
\left\{w_k\phi(k)+2\sum_{j<k}w_j\phi(j)\right\}.
\label{eq:variance-prefix-collapse}
\end{align}
This is \eqref{eq:exact-class-variance}.
\end{proof}

\begin{remark}[Exact formula versus numerical evaluation]
\label{rem:exact-versus-binary64}
Equations \eqref{eq:exact-class-mean} and
\eqref{eq:exact-class-variance} are exact finite identities.  Decimal values
reported later are binary64 evaluations with explicit forward-error bounds;
they are not directed-rounding interval evaluations.
\end{remark}

\subsection{The exact order-dependent term}
\label{subsec:order-dependent-term}

The elementary square increment is
\begin{equation}
T_j^2-T_{j-1}^2=1+2\xi_jT_{j-1}.
\label{eq:support-square-increment}
\end{equation}
Define the future plateau weight
\begin{equation}
W_b=\sum_{j=b}^{m}w_j.
\label{eq:future-plateau-weight}
\end{equation}

\begin{proposition}[Diagonal term plus signed restoration]
\label{prop:diagonal-restoration}
Every member of \(\Om_N^{(1)}\) satisfies
\begin{align}
L_N\QF_{N_0,N}(w)
&=\Sigma_1(N_0,N)+2\mathcal J_N(w),
\label{eq:diagonal-restoration}\\
\mathcal J_N(w)
&=\sum_{b=2}^{m}W_b\xi_bT_{b-1}.
\label{eq:restoration-functional}
\end{align}
The first term is a mask invariant.  Consequently the complete class is
ordered by the single scalar \(\mathcal J_N\).
\end{proposition}

\begin{proof}
Insert \eqref{eq:support-square-increment} into
\(T_j^2=\sum_{b\leq j}(T_b^2-T_{b-1}^2)\), multiply by \(w_j\), and interchange
the two finite sums.  The diagonal contribution is
\(\sum_jjw_j=\Sigma_1\).
\end{proof}

Exchangeability also gives
\begin{equation}
\E_{\pi}(\xi_bT_{b-1})=(b-1)\mu_2,
\label{eq:expected-restoration-step}
\end{equation}
and hence
\begin{equation}
\E_{\pi}\mathcal J_N
=\mu_2\sum_{b=2}^{m}(b-1)W_b.
\label{eq:expected-restoration-functional}
\end{equation}
For the nearly balanced classes in the data, \(\mu_2\) is close to
\(-1/(m-1)\).  The much smaller observed value of \(\QF\) therefore means
that the ordered products \(\xi_bT_{b-1}\), with their future weights, are
substantially more negative than their exchangeable-class average.

\subsection{Exact change under a transposition}
\label{subsec:transposition-calculus}

Let \(a<b\) be support indices carrying opposite signs.  Swapping them
preserves the mask, sign counts, and endpoint.

\begin{proposition}[Opposite-sign transposition formula]
\label{prop:transposition-formula}
If \(\xi_a=+1\) and \(\xi_b=-1\), then
\begin{equation}
L_N\{\QF(\pi\circ\tau_{a,b})-\QF(\pi)\}
=4\sum_{j=a}^{b-1}w_j(1-T_j).
\label{eq:plus-minus-transposition}
\end{equation}
If \(\xi_a=-1\) and \(\xi_b=+1\), then
\begin{equation}
L_N\{\QF(\pi\circ\tau_{a,b})-\QF(\pi)\}
=4\sum_{j=a}^{b-1}w_j(1+T_j).
\label{eq:minus-plus-transposition}
\end{equation}
\end{proposition}

\begin{proof}
In the first case the path decreases by \(2\) on \(a\leq j<b\) and is
unchanged elsewhere, so
\((T_j-2)^2-T_j^2=4(1-T_j)\).  The second case is analogous.
\end{proof}

With
\begin{equation}
\overline T_w[a,b)
=\frac{\sum_{j=a}^{b-1}w_jT_j}
{\sum_{j=a}^{b-1}w_j},
\label{eq:weighted-interval-mean}
\end{equation}
a \(+,-\) transposition decreases the quadratic mass exactly when
\(\overline T_w[a,b)>1\), and a \(-,+\) transposition decreases it exactly
when \(\overline T_w[a,b)<-1\).  Thus a local minimum under all admissible
transpositions is characterized by these two orientation-sensitive families
of inequalities.  They are not equivalent to a single condition
\(|\overline T_w[a,b)|\leq1\) over every interval, because not every interval
has both endpoint orientations.

The lower set determined by the observed word is therefore
\begin{equation}
\cA_N
=\{\pi\in\Om_N^{(1)}:\QF_{N_0,N}(\pi)
\leq\QF_{N_0,N}(\mu)\}
=\{\pi:\mathcal J_N(\pi)\leq\mathcal J_N(\mu)\}.
\label{eq:lower-mass-set}
\end{equation}
Finding the best single transposition takes \(O(m)\) time using running
extrema of weighted prefixes.  Counting every improving pair is a different
problem and generally requires more than that single pass.

\subsection{Occupation interpretation}
\label{subsec:occupation-interpretation}

Put \(d=1/\zeta(2)=6/\pi^2\), the standard squarefree density
\cite{MontgomeryVaughan2007}.  Since the \(j\)-th support plateau is centered
near \(n=j/d\) and has average integer length \(1/d\), one expects
\begin{equation}
w_j\sim\frac{d}{j^2}.
\label{eq:plateau-weight-asymptotic}
\end{equation}
This gives the useful approximation
\begin{equation}
\QF_{N_0,N}(w)
\approx d\left\langle\frac{T_j^2}{j}\right\rangle_{\log j}.
\label{eq:logarithmic-occupation-approximation}
\end{equation}
The exact probability measure behind this statement is
\begin{equation}
\nu_N(\{j\})=\frac{jw_j}{D_N},\qquad
D_N=\sum_{j=1}^m jw_j=\Sigma_1(N_0,N).
\label{eq:occupation-measure}
\end{equation}
Its normalized quadratic occupation level is
\begin{equation}
A_N^{\mathrm{occ}}
=\int\frac{T_j^2}{j}\,d\nu_N(j)
=\frac{L_N\QF_{N_0,N}(w)}{\Sigma_1(N_0,N)}.
\label{eq:exact-occupation-level}
\end{equation}
For every \(\lambda>0\), Markov's inequality gives the exact consequence
\begin{equation}
\nu_N\left\{j:\frac{T_j^2}{j}>\lambda\right\}
\leq\frac{A_N^{\mathrm{occ}}}{\lambda}.
\label{eq:occupation-markov}
\end{equation}
At the full cutoff, the density-normalized approximation
\(a_N=\zeta(2)\QF(\mu)=0.048252\) implies, at the descriptive level, that the
path lies below \(\sqrt j\) for at least \(95.2\) percent of logarithmic
support time and below \(2\sqrt j\) for at least \(98.8\) percent.  The exact
percentages use \(A_N^{\mathrm{occ}}\) rather than \(a_N\); the difference is
the finite mask correction in \eqref{eq:occupation-measure}.

\subsection{Support-time coefficients and a feedback heuristic}
\label{subsec:support-time-coefficients}

For an interval \(u<j\leq v\), define
\begin{equation}
c_{u,v}=\frac{1}{v-u}\sum_{j=u+1}^{v}\frac{T_j^2}{j}.
\label{eq:local-support-coefficient}
\end{equation}
Under an explicitly specified ensemble, one may also define
\begin{equation}
c(K)=\frac{\Var(T_K)}{K}.
\label{eq:variance-scale-coefficient}
\end{equation}
These are distinct quantities.  The algebraic normalization
\begin{equation}
a_N=\zeta(2)\QF_{N_0,N}(\mu)
\label{eq:density-normalized-quadratic-mass}
\end{equation}
is not an independent measurement of either coefficient.

\begin{heuristic}[Critical restoring drift]
\label{heur:critical-restoring-drift}
This is a heuristic argument.  Suppose a process obeys
\begin{equation}
\E(\xi_{j+1}\mid T_j)\approx-\frac{2a}{j}T_j.
\label{eq:critical-restoring-drift}
\end{equation}
Writing \(V_j=\E T_j^2\) gives
\begin{equation}
V_{j+1}-V_j\approx1-\frac{4a}{j}V_j.
\label{eq:variance-feedback-recursion}
\end{equation}
A solution \(V_j\sim cj\) then satisfies
\begin{equation}
c=\frac{1}{1+4a},\qquad a=\frac{c^{-1}-1}{4}.
\label{eq:feedback-coefficient-relation}
\end{equation}
The power \(j^{-1}\) is critical: a drift \(j^{-\gamma}\) formally gives
\(V_j/j\to1\) for \(\gamma>1\), \(V_j/j\to0\) for \(\gamma<1\), and permits a
strict constant in \((0,1)\) only for \(\gamma=1\).  This model can explain
how a cumulative suppression coexists with a null nearest-neighbour
persistence statistic, but \eqref{eq:feedback-coefficient-relation} is not an
identity for the M\"obius word.
\end{heuristic}

The scale measurements in Table \ref{tab:support-coefficients} fluctuate
between approximately \(0.036\) and \(0.068\), without monotone convergence.
A sinusoidal fit in \(\log_2K\) is reported in Appendix
\ref{app:numerical-results}; its out-of-sample performance does not establish
a predictive law.

\section{Excursion degeneracy and exact return counts}
\label{sec:excursion-geometry}

\subsection{Reflection orbit}
\label{subsec:reflection-orbit}

Let
\begin{equation}
0=\tau_0<\tau_1<\cdots<\tau_K
\label{eq:completed-return-times}
\end{equation}
be the successive completed zero-to-zero excursions of the support-time path.
Write the \(r\)-th completed excursion as
\begin{equation}
T|_{[\tau_{r-1},\tau_r]}=\varepsilon_rH_r,\qquad
\varepsilon_r\in\{-1,+1\},\qquad H_r\geq0.
\label{eq:excursion-decomposition}
\end{equation}

\begin{proposition}[Pointwise absolute-path orbit]
\label{prop:reflection-orbit}
Reflecting any subset of the \(K\) completed excursions preserves the zero
mask, the two sign counts, the endpoint, and every value \(R_w(n)^2\).
Moreover, the resulting orbit exhausts the pointwise absolute-path class:
\begin{equation}
\cO_N(\mu)=\Om_N^{(2)},\qquad
|\cO_N(\mu)|=|\Om_N^{(2)}|=2^K,\qquad
\QF_{N_0,N}(w)=\QF_{N_0,N}(\mu)
\quad(w\in\cO_N(\mu)).
\label{eq:reflection-orbit-size}
\end{equation}
If \(M(N)\neq0\), the terminal incomplete excursion after the last return is
not reflected, because its reflection would change the fixed endpoint.
\end{proposition}

\begin{proof}
A completed excursion contains equally many positive and negative steps.
Reflection changes every sign in that segment, preserves those counts, and
replaces the path by its negative only within the segment.  Independent
choices therefore give \(2^K\) distinct words in \(\Om_N^{(2)}\).

Conversely, take \(w\in\Om_N^{(2)}\).  Between two consecutive zeros
\(\tau_{r-1}\) and \(\tau_r\), the reference path has no interior zero.
Because both paths have unit increments in support time and the same absolute
value at every support index, the sign relating them cannot change inside
that interval.  Hence the restriction of \(w\) to the excursion is either
the original sign block or its complete reflection.  These alternatives are
independent across completed excursions.  On a terminal incomplete segment,
the fixed endpoint \(T_m=S\neq0\) forces the original orientation.  Thus
there are no further elements of \(\Om_N^{(2)}\), proving exhaustivity and
the asserted equality.  When \(S=0\), the terminal segment is itself the
last completed excursion.
\end{proof}

\begin{remark}[Consequence for selection]
\label{rem:reflection-selection-limit}
The reflection orbit describes degeneracy after the absolute path has already
been fixed.  It does not explain why the observed absolute path has small
quadratic mass.  It proves instead that no functional depending only on
\(R_w(n)^2\) can select a unique orientation.
\end{remark}

\subsection{Expected number of returns}
\label{subsec:expected-return-count}

For \(w\in\Om_N^{(1)}\), define
\begin{equation}
Z_N(w)=\#\{1\leq j\leq m:T_j(w)=0\}.
\label{eq:return-count}
\end{equation}

\begin{proposition}[Exact first moment of the return count]
\label{prop:exact-return-mean}
Under the uniform fixed-count measure,
\begin{equation}
\E_{\pi}Z_N
=\sum_{k=1}^{\min(p,q)}
\frac{\binom pk\binom qk}{\binom m{2k}}.
\label{eq:exact-return-mean}
\end{equation}
\end{proposition}

\begin{proof}
A return is possible only at time \(2k\).  Among the
\(\binom m{2k}\) equally likely subsets of sampled signs, precisely
\(\binom pk\binom qk\) contain \(k\) signs of each type.  Sum the indicators
of the return events.
\end{proof}

At \(N=5\cdot10^6\), the exact and observed values reported in Appendix
\ref{app:numerical-results} are
\begin{equation}
\E_{\pi}Z_N=1623.0496,\qquad
Z_N(\mu)=6016,\qquad
\frac{Z_N(\mu)}{\E_{\pi}Z_N}=3.7066.
\label{eq:five-million-return-excess}
\end{equation}
The same path has longest completed-excursion fraction \(0.11737\), compared
with a reported bridge-sample mean \(0.5301\).  These are related
descriptions of a path spending unusually little logarithmic mass far from
zero; they are not independent events whose probabilities may be
multiplied.

\section{Finite placement of the M\"obius word}
\label{sec:finite-placement}

Define the class-relative ratio and standardized displacement
\begin{equation}
\rho_Q(N)=\frac{\QF_{N_0,N}(\mu)}{\overline{\QF}_N},
\qquad
Z_Q(N)=\frac{\QF_{N_0,N}(\mu)-\overline{\QF}_N}
{\sqrt{\Var_{\pi}(\QF_{N_0,N})}}.
\label{eq:ratio-and-standardized-displacement}
\end{equation}

\begin{proposition}[Finite separation]
\label{prop:finite-separation}
Let the input be the signed-byte file with 455,052,511 entries and SHA-256
digest
\begin{equation}
\texttt{0e551f47d924ef6f3e348e131fdf53e0c677ad5e5471ef9d6249b3fd662e2ea4}.
\label{eq:data-sha256}
\end{equation}
Under the binary64 and execution assumptions specified in Appendix
\ref{app:reproducibility}, the fourteen cutoffs in Table
\ref{tab:full-fixed-class-table} satisfy
\begin{equation}
\QF_{1000,N}(\mu)<\overline{\QF}_N.
\label{eq:finite-separation-inequality}
\end{equation}
At the complete cutoff,
\begin{equation}
\QF_{1000,455052511}(\mu)=0.029334,\quad
\overline{\QF}_{455052511}=0.56628,\quad
\rho_Q(455052511)=0.0518.
\label{eq:complete-cutoff-values}
\end{equation}
At \(N=5\cdot10^6\), the higher-precision reported values are
\begin{equation}
\QF_{1000,N}(\mu)=0.0295150019629,\quad
\overline{\QF}_N=0.548481916755,\quad
\rho_Q(N)=0.0538121696655.
\label{eq:five-million-values}
\end{equation}
\end{proposition}

\begin{proof}
The class mean is evaluated from Theorem \ref{thm:exact-class-mean}; the
observed value is a positive weighted sum.  The streaming implementation uses
blocked summation, compensation between blocks, and an a-priori forward-error
propagation under the assumptions in Appendix \ref{app:reproducibility}.
For each cutoff the upper endpoint of the propagated interval for
\(\QF(\mu)-\overline{\QF}_N\) is negative.  The smallest ratio of the
absolute separation to the propagated bound is \(5.75\cdot10^{10}\).
Division gives \eqref{eq:complete-cutoff-values} and
\eqref{eq:five-million-values}.  The proposition is conditional on the
identified file and execution model; the exact symbolic identities do not
depend on either.
\end{proof}

\begin{remark}[What has been proved]
\label{rem:finite-not-asymptotic}
Proposition \ref{prop:finite-separation} proves comparison with the exact
class mean at fourteen finite cutoffs.  It does not prove that
\(\rho_Q(N)\) is always near \(0.05\), that the ratio is monotone, or that it
has a limit.  Indeed the first tabulated value is \(0.082224\), and the later
values exhibit small nonmonotone fluctuations.
\end{remark}

At the complete cutoff,
\begin{equation}
m=276{,}638{,}778,\qquad S=-5414,\qquad
\max_{n\leq N}|M(n)|=6695.
\label{eq:complete-cutoff-path-data}
\end{equation}
Since \(\sqrt m=16632\) to the displayed precision, the observed maximum is
\begin{equation}
\frac{6695}{\sqrt m}=0.402,\qquad
\frac{6695}{0.868731\sqrt m}=0.463.
\label{eq:observed-maximum-relative-scale}
\end{equation}
The second denominator is the limiting mean maximum of a nearly balanced
finite bridge.  This is a finite comparison; it is not a law-of-the-iterated-
logarithm statement.

\subsection{Lower-tail sampling}
\label{subsec:lower-tail-sampling}

The exact lower-tail fraction is
\begin{equation}
P_Q(N)=\frac{\#\{\pi\in\Om_N^{(1)}:
\QF_{N_0,N}(\pi)\leq\QF_{N_0,N}(\mu)\}}
{|\Om_N^{(1)}|}.
\label{eq:exact-lower-tail-fraction}
\end{equation}
No closed form for \eqref{eq:exact-lower-tail-fraction} is established here.
Uniform fixed-count sampling reported 412 hits in 300,000 draws at
\(N=10^4\), 15 hits in 150,000 draws at \(N=3\cdot10^4\), and zero hits at
each of the four larger sampled cutoffs.  Full results and exact zero-hit
bounds appear in Table \ref{tab:monte-carlo-ladder}.

\begin{experiment}[Failure of a two-moment tail approximation]
\label{exp:two-moment-rejection}
A scaled chi-square law fitted to the exact mean and variance predicts,
across the six cutoffs, expected hit counts
\begin{equation}
(20772,\ 6113.7,\ 2218.2,\ 477.88,\ 61.85,\ 23.242),
\label{eq:chi-square-expected-hits}
\end{equation}
whereas the observed counts are
\begin{equation}
(412,\ 15,\ 0,\ 0,\ 0,\ 0).
\label{eq:observed-hit-counts}
\end{equation}
At \(N=5\cdot10^6\), the approximation gives
\begin{equation}
P_Q^{(\chi^2)}=0.0023242189,\qquad
\Prob_{\chi^2}(H=0)=7.8397\cdot10^{-11}.
\label{eq:chi-square-zero-probability}
\end{equation}
The binomial probability of an outcome at least this extreme is below
\(7.84\cdot10^{-11}\) in every row.  Thus exact mean and variance do not
determine the lower tail of this quadratic form.
\end{experiment}

\begin{warning}[Mean separation is not a percentile]
\label{warn:mean-not-percentile}
The ratio \(\rho_Q\) compares the observed value with the exact mean; it is
not \(P_Q\).  Likewise, zero hits in \(R\) draws yield the exact one-sided
95-percent bound \(P_Q<1-0.05^{1/R}\), not an estimate of a probability of
order \(10^{-9}\).  At \(N=5\cdot10^6\), the class size is approximately
\(10^{915017.333}\), so exhaustive enumeration is impossible.
\end{warning}

\subsection{Continuum diagnostics}
\label{subsec:continuum-diagnostics}

Put
\begin{equation}
d=\frac{1}{\zeta(2)},\qquad
y_N^2=\frac{S^2}{m},\qquad x_0=e^{-L_N}.
\label{eq:continuum-parameters}
\end{equation}

\begin{heuristic}[Finite-span Gaussian bridge]
\label{heur:finite-span-gaussian-bridge}
This is a heuristic continuum replacement.  Replacing the squarefree mask by
density \(d\) and the hypergeometric bridge by a Gaussian bridge gives
\begin{align}
\overline{\QF}_N^{\mathrm{cont}}
&=d\left[1-(1-y_N^2)\frac{1-e^{-L_N}}{L_N}\right],
\label{eq:continuum-mean}\\
(\sigma_N^{\mathrm{cont}})^2
&=\frac{d^2}{L_N^2}\left[
4L_N-10+8(L_N+1)e^{-L_N}+2e^{-2L_N}\right.
\notag\\
&\hspace{35mm}\left.
+4y_N^2\{1-e^{-2L_N}-2L_Ne^{-L_N}\}\right].
\label{eq:continuum-variance}
\end{align}
Dropping only the exponentially small term in the mean gives
\begin{equation}
\overline{\QF}_N^{\mathrm{principal}}
=d\left[1-\frac{1-y_N^2}{L_N}\right].
\label{eq:principal-continuum-mean}
\end{equation}
When \(y_N^2=o(1)\),
\begin{equation}
\overline{\QF}_N^{\mathrm{cont}}
=d\left(1-\frac1{L_N}\right)+o(L_N^{-1}),\qquad
(\sigma_N^{\mathrm{cont}})^2
=\frac{4d^2(L_N-5/2)}{L_N^2}+o(L_N^{-2}).
\label{eq:continuum-asymptotics}
\end{equation}
\end{heuristic}

At the complete cutoff, \eqref{eq:principal-continuum-mean} gives
\(0.5662087354\), while \eqref{eq:continuum-mean} gives
\(0.5662088271\).  The exact finite-class mean is \(0.56628\).  The exact and
continuum standard deviations are \(0.30436\) and \(0.30433\), respectively.
These close numerical agreements validate the approximation at the measured
scale but do not prove its asymptotic error term.

\begin{heuristic}[Ornstein--Uhlenbeck lower-tail rate]
\label{heur:ou-rate}
The Lamperti transform
\begin{equation}
V(t)=e^{t/2}W(e^{-t}),\qquad
\Cov(V(s),V(t))=e^{-|t-s|/2}
\label{eq:lamperti-transform}
\end{equation}
is a stationary Ornstein--Uhlenbeck process with unit variance.  The
Feynman--Kac calculation within the Donsker--Varadhan framework
\cite{DonskerVaradhan1975} gives
\begin{align}
\Lambda(\lambda)
&=\frac12\left(\frac12-\sqrt{\frac14-2\lambda}\right),
\qquad\lambda<\frac18,
\label{eq:ou-cumulant}\\
I(a)
&=\sup_{\lambda<1/8}\{\lambda a-\Lambda(\lambda)\}
=\frac{(a-1)^2}{8a}.
\label{eq:ou-rate-function}
\end{align}
Transferring this rate to the finite fixed-mask bridges requires an unproved
large-deviation invariance principle and a boundary correction.  The
resulting prediction is
\begin{equation}
\log P_Q(N)\approx-L_NI(a_N)+\log\kappa(a_N,L_N),
\qquad a_N=\zeta(2)\QF_{N_0,N}(\mu),
\label{eq:ou-lower-tail-prediction}
\end{equation}
where the prefactor \(\kappa\) is not determined.
\end{heuristic}

At \(N=5\cdot10^6\), the leading rate gives
\begin{equation}
a_N=0.0485502,\qquad I(a_N)=2.3307,\qquad
\exp\{-L_NI(a_N)\}=2.39\cdot10^{-9}.
\label{eq:ou-five-million-prediction}
\end{equation}
At the complete cutoff the corresponding values are
\begin{equation}
a_N=0.048252,\qquad I(a_N)=2.3466,\qquad
\exp\{-L_NI(a_N)\}=5.28\cdot10^{-14}.
\label{eq:ou-complete-prediction}
\end{equation}
These are heuristic extrapolations far below the resolution of the uniform
sampler.  The straight line through the two cutoffs having positive hit
counts has slope \(-2.38467\); with only two points it carries no residual
regression degrees of freedom.  A descriptive six-cutoff binomial likelihood
fit gives exponent \(2.45306\) and profile interval
\([2.04052,2.94753]\).

\section{Classical analytic meaning of the quadratic scale}
\label{sec:analytic-meaning}

\subsection{Full-prefix and truncated normalizations}
\label{subsec:normalization-comparison}

The classical full-prefix quantity is
\begin{equation}
\mathfrak Q(N)
=\frac{1}{\log N}\sum_{n\leq N}\frac{M(n)^2}{n^2}.
\label{eq:full-prefix-quadratic-mass}
\end{equation}
For fixed \(N_0\), set
\begin{equation}
K_{N_0}=\sum_{n\leq N_0}\frac{M(n)^2}{n^2}.
\label{eq:prefix-constant}
\end{equation}
Then the exact conversion is
\begin{equation}
\mathfrak Q(N)
=\frac{K_{N_0}+L_N\QF_{N_0,N}(\mu)}{\log N}.
\label{eq:normalization-conversion}
\end{equation}
The two normalizations have the same limit whenever either limit exists, but
they can be far apart at finite \(N\).  With \(N_0=1000\),
\begin{equation}
K_{1000}=2.068991,\qquad
\mathfrak Q(455052511)=0.122952,
\label{eq:prefix-numerical-values}
\end{equation}
whereas the truncated value is \(0.029334\).  Thus the observed flatness of
the latter is a property of the fixed lower cutoff, not evidence that the
full-prefix normalization has numerically converged.

\subsection{Weak Mertens and RH}
\label{subsec:weak-mertens-rh}

For integer \(N\),
\begin{equation}
\int_1^{N+1}\left(\frac{M(x)}{x}\right)^2\,dx
=\sum_{n=1}^{N}\frac{M(n)^2}{n(n+1)}
=\sum_{n=1}^{N}\frac{n}{n+1}\frac{M(n)^2}{n^2},
\qquad \frac12\leq\frac{n}{n+1}<1.
\label{eq:discrete-continuous-comparison}
\end{equation}
Hence boundedness of \(\mathfrak Q(N)\) is equivalent, up to absolute
comparison constants, to the weak Mertens conjecture
\begin{equation}
\int_1^X\left(\frac{M(x)}{x}\right)^2\,dx\ll\log X.
\label{eq:weak-mertens-conjecture}
\end{equation}
This conjecture implies RH and simplicity of all nontrivial zeros; see
\cite{Titchmarsh1986,OdlyzkoTeRiele1985,Ng2004}.  RH alone is not known to
imply \eqref{eq:weak-mertens-conjecture}.

There is also a direct positive reformulation.  Define
\begin{align}
D(N)&=\sum_{k\leq N}\frac{\mu(k)^2}{k},
\label{eq:diagonal-bilinear-sum}\\
C(N)&=\sum_{k\leq N}\frac{\mu(k)M(k-1)}{k},
\label{eq:cross-bilinear-sum}\\
B(N)&=D(N)+2C(N).
\label{eq:positive-bilinear-functional}
\end{align}
Discrete summation by parts gives
\begin{equation}
B(N)=\frac{M(N)^2}{N}
+\sum_{k<N}\frac{M(k)^2}{k(k+1)}\geq0.
\label{eq:positive-bilinear-representation}
\end{equation}
It follows directly from the classical Mertens criterion that
\begin{equation}
\mathrm{RH}\quad\Longleftrightarrow\quad
B(N)\ll_{\eps}N^{\eps}\quad\text{for every }\eps>0.
\label{eq:positive-rh-criterion}
\end{equation}
This is a reformulation, not a new criterion for RH.

\begin{proposition}[Conditional meaning of a bounded class ratio]
\label{prop:conditional-ratio-meaning}
Assume
\begin{equation}
\frac{M(N)^2}{Q(N)\log(N/N_0)}\longrightarrow0,
\label{eq:endpoint-negligibility}
\end{equation}
equivalently \(M(N)=o(\sqrt{N\log N})\).  Then
\begin{equation}
\overline{\QF}_N\longrightarrow d=\frac1{\zeta(2)}
\label{eq:null-mean-limit}
\end{equation}
and
\begin{equation}
\sup_N\rho_Q(N)<\infty
\quad\Longleftrightarrow\quad
\QF_{N_0,N}(\mu)=O(1)
\quad\Longleftrightarrow\quad
\mathfrak Q(N)=O(1).
\label{eq:ratio-weak-mertens-equivalence}
\end{equation}
\end{proposition}

\begin{proof}
The elementary estimate \(Q(n)=dn+O(\sqrt n)\) gives
\(\Sigma_1=dL_N+O(1)\).  In
\eqref{eq:exact-class-mean-alternative}, the bridge correction is
\(O(L_N^{-1})\), while the endpoint contribution is
\(O(M(N)^2/[Q(N)L_N])\).  Equation \eqref{eq:endpoint-negligibility} proves
\eqref{eq:null-mean-limit}.  The first equivalence in
\eqref{eq:ratio-weak-mertens-equivalence} follows because the mean is then
bounded above and away from zero; the second follows from
\eqref{eq:normalization-conversion}.
\end{proof}

\begin{warning}[The endpoint hypothesis is not an RH consequence]
\label{warn:endpoint-hypothesis}
The standard consequence of RH,
\(M(N)=O_{\eps}(N^{1/2+\eps})\), does not imply
\(M(N)=o(\sqrt{N\log N})\).  Assuming the stronger
\(M(N)=O(\sqrt N)\) would make boundedness of the quadratic mass immediate.
Thus Proposition \ref{prop:conditional-ratio-meaning} cannot be invoked from
RH alone.
\end{warning}

\begin{remark}[The status of an asymptotic five-percent claim]
\label{rem:always-five-percent}
A statement such as
\begin{equation}
\rho_Q(N)\leq0.05+o(1)
\label{eq:five-percent-asymptotic-claim}
\end{equation}
would be a quantitative strengthening of bounded logarithmic second moment,
not an elementary consequence of the finite model.  Even the weaker
assertion \(\sup_N\rho_Q(N)<\infty\) requires, under
\eqref{eq:endpoint-negligibility}, the weak Mertens conjecture.  The finite
data do not prove \eqref{eq:five-percent-asymptotic-claim}; they only show
ratios near five percent over the stated cutoffs.
\end{remark}

\subsection{Conditional spectral constant}
\label{subsec:spectral-constant}

Under RH together with the negative-moment hypothesis
\(J_{-1}(T)\ll T\), Ng proves the logarithmic second-moment limit
\cite{Ng2004}
\begin{align}
\frac1Y\int_0^Y e^{-y}M(e^y)^2\,dy
&\longrightarrow\beta_M,
\label{eq:ng-second-moment}\\
\beta_M
&=2\sum_{\gamma>0}\frac1{|\rho\zeta'(\rho)|^2}.
\label{eq:ng-spectral-constant}
\end{align}
A separate reported computation gives
\begin{align}
2\sum_{k=1}^{400}\frac1{|\rho_k\zeta'(\rho_k)|^2}
&=0.0287674,
\label{eq:four-hundred-zero-sum}\\
\widehat R_{400}
=\frac{2\widehat c}{T_{400}}
&=0.0002715,\qquad
\widehat c=0.0923,\quad T_{400}=679.74,
\label{eq:estimated-spectral-tail}\\
\widehat\beta_M
&=0.0290389.
\label{eq:tail-corrected-spectral-constant}
\end{align}
The second term is a tail ansatz, not a computed sum.  It estimates the
positive-ordinate density of \(|\zeta'(\rho)|^{-2}\) from the last quarter of
the 400 ordinates.

Gonek's conjecture is
\begin{equation}
J_{-1}(T)
:=\sum_{0<\gamma\leq T}\frac1{|\zeta'(\rho)|^2}
\sim\frac{3}{\pi^3}T,
\label{eq:gonek-conjecture}
\end{equation}
where \(3/\pi^3\), not \(6/\pi^3\), is the constant for positive ordinates
\cite{Gonek1989,Ng2004}.  Substitution into the same tail ansatz gives
\begin{equation}
\frac{2(3/\pi^3)}{T_{400}}=0.0002847,\qquad
0.0287674+\frac{2(3/\pi^3)}{T_{400}}=0.0290521.
\label{eq:gonek-tail-comparison}
\end{equation}
The empirical \(\widehat c=0.0923\) is \(4.6\) percent below
\(3/\pi^3=0.0967546\ldots\), so there is no factor-two discrepancy.

Under the hypotheses needed for \eqref{eq:ng-second-moment} and for the null
mean limit,
\begin{equation}
\rho_Q(N)\longrightarrow
\rho_{\infty}=\zeta(2)\beta_M.
\label{eq:conditional-ratio-limit}
\end{equation}
The tail-corrected 400-zero value gives
\begin{equation}
\widehat\rho_{\infty}=0.0477671,\qquad
I(\widehat\rho_{\infty})=2.37284.
\label{eq:conditional-spectral-ratio}
\end{equation}
This explains why a ratio near five percent is plausible.  It also answers
the novelty question sharply: the approximate asymptotic percentage is a
direct conditional combination of a classical spectral constant with the
fixed-class mean \(d\); the finite exact denominator and its comparison with
the observed word are the model-specific contribution.

The reported spectral mass is concentrated at low ordinates:
\begin{equation}
\text{first zero: }54.7\%,\qquad
\text{first five: }80.3\%,\qquad
\text{first one hundred: }97.4\%.
\label{eq:low-zero-concentration}
\end{equation}
Thus the small quadratic scale is principally a low-frequency property.
The same separate computation reports a zero-parameter reconstruction of the
fourteen truncated values with RMS error \(1.5\cdot10^{-6}\).  These spectral
numbers require the separate reproducibility package specified in Appendix
\ref{app:reproducibility}; they are not inputs to Proposition
\ref{prop:finite-separation}.

\begin{warning}[Quadratic mass does not control remote maxima]
\label{warn:quadratic-mass-not-maximum}
The constant in \eqref{eq:ng-spectral-constant} is an \(\ell^2\)-type norm of
zero coefficients.  It does not bound their \(\ell^1\)-norm, force phase
alignment, or control \(\sup_x|M(x)|/\sqrt x\).  Finite flatness of
\(\QF_{1000,N}\) therefore gives no evidence that the normalized Mertens
maximum remains globally bounded.
\end{warning}

\section{Structural experiments beyond quadratic mass}
\label{sec:structural-experiments}

The tests in this section ask whether the observed order contains information
not already summarized by the quadratic functional.  Their status is
experimental, and their numerical specifications are recorded in Appendices
\ref{app:numerical-results} and \ref{app:reproducibility}.

\subsection{Excursion orientation and shape}
\label{subsec:orientation-shape-tests}

For each completed excursion, let \(\varepsilon_r\) be its orientation,
\(\ell_r\) its support-time length, and
\begin{equation}
s_r=\sum_{n\ \mathrm{in\ excursion}\ r}\frac1n
\label{eq:excursion-log-duration}
\end{equation}
its logarithmic integer duration.  The tested orientation statistics were
\begin{align}
\Theta&=\frac{\sum_rs_r\varepsilon_r}{\sum_rs_r},
\label{eq:orientation-bias-statistic}\\
P_1&=\frac1{K-1}\sum_{r<K}\varepsilon_r\varepsilon_{r+1},
\label{eq:orientation-persistence-statistic}\\
A_{\ell}&=\Corr(\varepsilon_r,\log\ell_r).
\label{eq:orientation-length-statistic}
\end{align}
Conditional randomization with the unsigned excursion skeleton fixed gave
\begin{equation}
p_{\Theta}=0.574,\qquad
p_{P_1}=0.271,\qquad
p_{A_{\ell}}=0.479.
\label{eq:orientation-test-results}
\end{equation}
No orientation bias, nearest-neighbour persistence, or sign--length
association was detected.

For unsigned excursion shapes, the tested coordinates were
\begin{align}
A_r^{\mathrm{sh}}
&=\frac1{\ell_r^2}\sum_tH_r(t)^2,
\label{eq:shape-area-statistic}\\
H_r^{\max}
&=\frac{\max_tH_r(t)}{\sqrt{\ell_r}},
\label{eq:shape-height-statistic}\\
\operatorname{HS}_r
&=\text{Horton--Strahler order}.
\label{eq:shape-tree-statistic}
\end{align}
Five hundred observed excursions and 100 conditioned Catalan replicates gave
\begin{equation}
p_A=0.505,\qquad p_H=0.743,\qquad p_{\mathrm{HS}}=0.485.
\label{eq:catalan-shape-results}
\end{equation}
The tested shape statistics therefore do not distinguish the observed sample
from the conditioned Catalan reference.  Catalan enumeration and uniform
Dyck paths are standard; see \cite{Stanley1999,FlajoletSedgewick2009}.

\subsection{Scale partition and slow-envelope controls}
\label{subsec:scale-controls}

The partition audit in Table \ref{tab:partition-audit} compares an unweighted
bin standard deviation with a Pearson-normalized statistic.  The former has a
two-sided value \(p=0.0066\) at 24 bins but not at 12 or 48 bins; the latter
is nonsignificant at all three resolutions.  There is therefore no
partition-robust result.

For sliding windows, the short-scale statistics at window sizes 256 and 512
have correlations
\begin{equation}
r_{256}=0.931,\qquad r_{512}=0.914,\qquad p<10^{-4}
\quad\text{in both cases}
\label{eq:slow-envelope-correlations}
\end{equation}
with the local trend \(M(x)/\sqrt x\).  At window size 1024, the reported
split-half variance ratios are
\begin{equation}
(0.248,\ p=0.111)\quad\text{and}\quad
(0.706,\ p=0.687).
\label{eq:split-half-results}
\end{equation}
The familywise result over seven scales from 256 through 2978 is
\begin{equation}
p_{\mathrm{global}}=0.0645.
\label{eq:scale-family-result}
\end{equation}
These controls identify the smaller-scale feature with the known slow
Mertens envelope and show that the larger-scale feature does not replicate
between halves.

\subsection{Conditional-choice analysis}
\label{subsec:conditional-choice-analysis}

For blocks of \(q\) nonzero entries, the feature set included second- and
third-order R\'enyi entropies, zero local time, self-intersection count, path
width, turn rate, the quadratic functional, transposition stability, and
their prescribed interactions.  Given a choice set \(\mathcal S_r\), the
conditional-logit model was
\begin{equation}
\Prob(\pi_j\mid\mathcal S_r)
=\frac{\exp(\beta^{\mathsf T}X_{rj})}
{\sum_{\ell}\exp(\beta^{\mathsf T}X_{r\ell})}.
\label{eq:conditional-logit-model}
\end{equation}
See \cite{McFadden1974} for conditional logit and
\cite{BenjaminiHochberg1995} for false-discovery control.

At \(q=256,512,1024\), the quadratic-only deviance tests gave
\begin{equation}
p_{\mathrm{ANOVA}}=(0.550,\ 0.904,\ 0.661),
\label{eq:local-anova-results}
\end{equation}
with held-out gains per set
\begin{equation}
(1.0\cdot10^{-6},\ -4.16\cdot10^{-4},\ 0.00245).
\label{eq:local-anova-test-gains}
\end{equation}
The complete feature hierarchy did not produce a robust held-out selector.
These results are distinct from the blind spectral experiment below.

\subsection{Blind recovery of logarithmic frequencies}
\label{subsec:blind-spectral-recovery}

Define the log-time signal
\begin{equation}
Y_{\theta}(t)=e^{-\theta t}M(\lfloor e^t\rfloor).
\label{eq:blind-log-signal}
\end{equation}
A hierarchical spectral procedure chose \(\theta\) and frequencies from the
Mertens path without reading zeta zeros or their derivatives.  The external
reference audit was performed only after fitting.  Before the permutation
null filter, the procedure matched 10 of the first 14 reference ordinates.
For the first four it reported
\begin{equation}
\begin{array}{c|rrrr}
\text{reference ordinate}&14.1347&21.0220&25.0109&30.4249\\ \hline
\text{blind estimate}&14.138&21.017&25.011&30.430
\end{array}.
\label{eq:first-four-blind-frequencies}
\end{equation}
The full fourteen-mode audit appears in Table
\ref{tab:blind-frequency-audit}.

The complete discovery, validation, and test pipeline was then rerun on 100
uniform fixed-mask permutations.  A mode-by-mode filter at
\(q=0.010\) retained
\begin{equation}
0\ \text{of the 11 pre-null candidate modes}.
\label{eq:blind-null-filter-result}
\end{equation}
With the add-one empirical convention, however, the smallest attainable
modewise value and the first Benjamini--Hochberg threshold were
\begin{equation}
\widehat p_{\min}=\frac1{100+1}=0.00990099\ldots,
\qquad
\frac{q}{11}=0.00090909\ldots.
\label{eq:blind-bh-resolution-obstruction}
\end{equation}
Thus the design had no resolution for a sparse rank-one discovery at the
declared \(q\).  At least
\begin{equation}
R\geq\left\lceil\frac{11}{0.010}\right\rceil-1=1099
\label{eq:minimum-placebos-for-bh}
\end{equation}
complete placebo reruns are required merely to make the first BH threshold
attainable.  A prespecified global maximum statistic would answer a different
familywise question and could use the 100 reruns more efficiently.

\begin{remark}[Interpretation of the blind spectrum]
\label{rem:blind-spectrum-interpretation}
The pre-null frequency recovery is real and should not be omitted, but it is
not a discovery of previously unknown zeros.  Classical explicit formulae
place zero ordinates in the log-frequency content of \(M(x)/\sqrt x\); a
successful blind fit is therefore an algorithmic recovery of known spectral
structure.  The observation that no candidate survived the 100-placebo
modewise filter is inconclusive, because
\eqref{eq:blind-bh-resolution-obstruction} makes that filter underpowered by
construction for a sparse discovery.  The run neither establishes a new
fixed-class spectral effect nor supplies strong evidence that no such effect
exists.  It is also not the ANOVA in
\eqref{eq:local-anova-results}: the two procedures answer different
questions.
\end{remark}

\section{A route toward genuine arithmetic selection}
\label{sec:selection-route}

\subsection{The exact selector already known}
\label{subsec:exact-arithmetic-selector}

The fixed-mask class discards the arithmetic relation that uniquely defines
the M\"obius function.  That relation can be restored explicitly.  For a word
\(w\) with \(w(1)=1\), define the divisor-convolution defect
\begin{equation}
\Delta_N(w)
=\sum_{n=2}^{N}\omega_n
\left|\sum_{d\mid n}w(d)\right|^2,
\qquad \omega_n>0.
\label{eq:divisor-convolution-defect}
\end{equation}

\begin{proposition}[Unique finite arithmetic selector]
\label{prop:unique-arithmetic-selector}
For every \(N\geq2\),
\begin{equation}
\Delta_N(w)=0
\quad\Longleftrightarrow\quad
w(n)=\mu(n)\ \text{for all }n\leq N.
\label{eq:unique-arithmetic-selection}
\end{equation}
\end{proposition}

\begin{proof}
The left side is equivalent to
\(\sum_{d\mid n}w(d)=0\) for every \(2\leq n\leq N\), together with
\(w(1)=1\).  These are the finite equations
\((w*\mathbf 1)(n)=\delta_{n1}\).  The Dirichlet inverse of the constant-one
function is unique and equals \(\mu\).  Equivalently, induction uses
\begin{equation}
w(n)=-\sum_{\substack{d\mid n\\d<n}}w(d)
\label{eq:mobius-recursion}
\end{equation}
to determine each value.
\end{proof}

This proposition gives a unique position to \(M\) in the strongest possible
sense, but it merely restates M\"obius inversion.  It contains no cancellation
estimate and therefore supplies no progress toward RH.

\subsection{Why the finite class cannot classify RH behavior}
\label{subsec:finite-class-rh-classification}

RH is a statement about one projectively consistent infinite arithmetic
sequence.  A member of \(\Om_N^{(1)}\) is only a finite word, and the uniform
classes for successive \(N\) do not define a canonical infinite extension.
It is therefore not mathematically meaningful to partition one finite class
into words that satisfy or violate RH without first prescribing an extension
law.

For an infinite sequence \(w\), one may define its growth exponent
\begin{equation}
h(w)=\limsup_{x\to\infty}
\frac{\log(1+|R_w(x)|)}{\log x}.
\label{eq:path-growth-exponent}
\end{equation}
For the M\"obius function, RH is equivalent to \(h(\mu)\leq1/2\).  The
stronger weak-Mertens property is
\begin{equation}
\sup_{N\geq2}\frac1{\log N}
\sum_{n\leq N}\frac{R_w(n)^2}{n^2}<\infty.
\label{eq:abstract-weak-second-moment}
\end{equation}
For arbitrary paths, \eqref{eq:abstract-weak-second-moment} does not control
the supremum because sparse spikes can be hidden in an average.  Its
implication for RH is specifically arithmetic and uses the Mellin transform
of \(M\) and \(1/\zeta\).

\subsection{The missing inequality}
\label{subsec:missing-coercive-inequality}

A nontrivial selection theorem would need more than
\(\Delta_N(\mu)=0\).  One viable target is a scale-consistent inequality that
links a relaxed divisor defect to a positive path functional, schematically
\begin{equation}
B_w(N)=\frac{R_w(N)^2}{N}
+\sum_{k<N}\frac{R_w(k)^2}{k(k+1)}
\label{eq:general-positive-path-functional}
\end{equation}
and
\begin{equation}
B_w(N)
\leq C_{\eps}N^{\eps}
+\mathcal R_N\!\left(\Delta_{\leq N}(w)\right),
\qquad
\mathcal R_N(0)=0,
\label{eq:desired-coercive-inequality}
\end{equation}
where \(B_w\) is the positive quantity analogous to
\eqref{eq:positive-bilinear-representation}.  Applied to \(w=\mu\), such a
bound would imply the RH criterion \eqref{eq:positive-rh-criterion}.  No
inequality of the required strength is proved here; obtaining it would be at
least as difficult as the target analytic estimate.

A statistically meaningful intermediate programme has three parts:
\begin{enumerate}[label=(\roman*)]
\item condition simultaneously on the mask, endpoint, and quadratic-mass
profile over many nested cutoffs, rather than on one scalar value;
\item test coordinates orthogonal to the absolute path, such as orientation
bits and cross-level nesting, under exact conditional nulls;
\item require projective consistency in \(N\) and a deterministic link from
any surviving coordinate to the divisor-convolution defect or to
\eqref{eq:positive-bilinear-representation}.
\end{enumerate}
The orientation and shape tests in Section \ref{sec:structural-experiments}
do not detect a departure on the available interval.  The blind spectrum recovers known
analytic structure but does not provide an independent combinatorial
selector.  Thus no current statistic uniquely distinguishes the M\"obius
path from all RH-compatible paths, let alone from paths satisfying a stronger
mean-square property.

\begin{question}[Projectively consistent residual selector]
\label{ques:projective-selector}
Does there exist a statistic \(\mathcal S_N(w)\), not determined by the mask,
endpoint, or multiscale quadratic-mass profile, whose conditional null is
explicit, whose value is consistent under increasing cutoffs, and for which a
deterministic inequality links \(\mathcal S_N(\mu)\) to
\eqref{eq:positive-rh-criterion}?
\end{question}

\section{Conclusion}
\label{sec:conclusion}

The model has a rigorous finite core and a limited analytic interpretation.
The fixed-mask, fixed-count orbit is exactly specified; its support-time path
is a hypergeometric bridge at the square-root critical scale.  The normalized
logarithmic quadratic path mass has exact finite mean and variance, an exact
order-dependent bilinear term, and an exact transposition calculus.  On the
identified data, the M\"obius word has \(5.18\) percent of its class mean at
the largest cutoff, with the inequality separated from zero by a large
forward-error margin.

For an expert in analytic number theory, the five-percent magnitude is not a
new asymptotic law: under Ng's hypotheses it is explained by
\(\zeta(2)\beta_M\), and the first few zeros dominate \(\beta_M\).  What the
model adds is a finite conditional baseline and a set of exact combinatorial
coordinates for asking how the sign order differs from that baseline.  The
reported lower tail is much thinner than a two-moment approximation, but its
exact probability remains unknown.

No tested excursion orientation, Catalan shape, scale partition, or local
conditional-choice statistic supplies a surviving selector beyond the
quadratic scale.  The blind spectral procedure recovers several known zero
frequencies without using them in the fit, but its permutation-corrected
mode filter retains none; with only 100 placebos, that filter is too coarse
for a sparse BH discovery and the outcome is therefore inconclusive.  The
only exact unique selector presently available
is the divisor-convolution identity, and connecting that identity to
square-root cancellation is the unresolved arithmetic problem rather than a
consequence of the finite permutation model.

\appendix

\section{Numerical tables}
\label{app:numerical-results}

This appendix collects the numerical outputs used in the text.  Exact
finite-class formulae are distinguished from sampled statistics in every
caption.

\subsection{Fixed-class quadratic comparison}
\label{subsec:numerical-fixed-class}

\begin{table}[htbp]
\centering
\scriptsize
\setlength{\tabcolsep}{3.2pt}
\caption{Observed quadratic path mass and exact fixed-mask, fixed-count class
moments, with \(N_0=1000\).  Decimal entries are binary64 evaluations of the
exact formulae.}
\label{tab:full-fixed-class-table}
\begin{tabular}{r r r r r r r}
\toprule
\(N\)&\(M(N)\)&\(\QF(\mu)\)&\(\overline{\QF}_N\)&\(\sigma_N\)&\(\rho_Q\)&\(Z_Q\)\\
\midrule
10,000&\(-23\)&0.032174&0.39129&0.37891&0.082224&\(-0.9478\)\\
30,000&18&0.030803&0.43847&0.39324&0.070251&\(-1.0367\)\\
100,000&\(-48\)&0.031084&0.48239&0.39632&0.064437&\(-1.1387\)\\
300,000&220&0.029369&0.53006&0.39912&0.055407&\(-1.2545\)\\
1,000,000&212&0.030100&0.52665&0.37332&0.057154&\(-1.3301\)\\
5,000,000&\(-709\)&0.029515&0.54848&0.35510&0.053812&\(-1.4614\)\\
50,000,000&\(-908\)&0.029733&0.55336&0.32471&0.053732&\(-1.6126\)\\
100,000,000&1,928&0.029641&0.55844&0.31817&0.053079&\(-1.6620\)\\
150,000,000&1,089&0.029536&0.55767&0.31333&0.052963&\(-1.6855\)\\
200,000,000&320&0.029772&0.55824&0.31038&0.053332&\(-1.7027\)\\
250,000,000&\(-982\)&0.029944&0.55941&0.30838&0.053529&\(-1.7169\)\\
300,000,000&\(-592\)&0.029662&0.55989&0.30663&0.052977&\(-1.7292\)\\
350,000,000&2,896&0.029508&0.56226&0.30578&0.052482&\(-1.7423\)\\
455,052,511&\(-5,414\)&0.029334&0.56628&0.30436&0.051800&\(-1.7642\)\\
\bottomrule
\end{tabular}
\end{table}

The complete-cutoff support contains \(276{,}638{,}778\) nonzero entries.
The corresponding squarefree density is \(0.60793\) to the displayed
precision.  All fourteen parity checks
\begin{equation}
M(N)\equiv Q(N)\pmod2
\label{eq:reported-parity-checks}
\end{equation}
passed.

\subsection{Uniform-permutation ladder}
\label{subsec:numerical-permutation-ladder}

\begin{table}[htbp]
\centering
\scriptsize
\caption{Uniform fixed-count sampling.  Positive-hit intervals are Wilson
intervals \cite{Wilson1927}; zero-hit entries are exact one-sided 95-percent
Clopper--Pearson bounds.}
\label{tab:monte-carlo-ladder}
\begin{tabular}{r r r r c r}
\toprule
\(N\)&draws&hits&\(\widehat P_Q\)&95-percent interval or bound&\(Z_s\)\\
\midrule
10,000&300,000&412&0.0013733&[0.0012470,0.0015124]&1.3216\\
30,000&150,000&15&0.0001000&[0.0000606,0.0001650]&0.9930\\
100,000&100,000&0&0&\(<2.9957\cdot10^{-5}\)&\(-1.6256\)\\
300,000&50,000&0&0&\(<5.9913\cdot10^{-5}\)&\(-0.7217\)\\
1,000,000&10,000&0&0&\(<2.9953\cdot10^{-4}\)&\(-1.9583\)\\
5,000,000&10,000&0&0&\(<2.9953\cdot10^{-4}\)&\(-1.2327\)\\
\bottomrule
\end{tabular}
\end{table}

Here \(Z_s=(s-\sigma_N)/\operatorname{SE}(s)\), with
\(\operatorname{SE}(s)\) estimated from the sample fourth moment.  The first
8,200 draws at \(N=5\cdot10^6\) have
\begin{equation}
\overline{\QF}_{\mathrm{sample}}=0.544194764161,\qquad
s_{\mathrm{sample}}=0.350499677391.
\label{eq:sample-mean-standard-deviation}
\end{equation}
These are compatible with the exact mean \(0.548481916755\) and exact
standard deviation \(0.35510\).  The value \(0.350195\) is the leading
\(L_N-5/2\) asymptotic standard deviation, not the exact finite-class value.

\subsection{Support-time coefficients}
\label{subsec:numerical-support-coefficients}

\begin{table}[htbp]
\centering
\small
\caption{Recorded scale coefficients.  The second and third comparison
columns arise from the feedback model and its scale adjustment; they are not
independent estimates of \(\zeta(2)\QF\).}
\label{tab:support-coefficients}
\begin{tabular}{r c c c}
\toprule
\(K\)&\(c_{\mathrm{observed}}\)&\(c_{a=4.8161}\)&\(c_{a,\mathrm{adjusted}}\)\\
\midrule
256&0.048729&0.047858&0.064982\\
512&0.067637&0.054504&0.066281\\
1024&0.036207&0.048046&0.065380\\
2048&0.049998&0.053648&0.037300\\
4096&0.054833&0.043456&0.050333\\
8192&0.044089&0.045434&0.046309\\
16,384&0.049748&0.050261&0.040051\\
32,768&0.038838&0.043234&0.044990\\
65,536&0.061641&0.050747&0.034831\\
131,072&0.039820&0.047357&0.062011\\
262,144&0.036670&0.047253&0.036656\\
524,288&0.058102&0.046882&0.035871\\
\bottomrule
\end{tabular}
\end{table}
The chosen reference
\begin{equation}
c_0=\frac{0.0300}{6/\pi^2}=\frac{\pi^2}{200}
=0.049348022005\ldots
\label{eq:chosen-feedback-reference}
\end{equation}
was fixed by the rounded target \(0.0300\), not estimated independently.
The first fitted log-scale mode was
\begin{equation}
g_1(t)=-0.000774584
-0.009187853\cos(2.541179t)
-0.006146902\sin(2.541179t),
\qquad t=\log_2K.
\label{eq:first-feedback-mode}
\end{equation}
Its period is \(2.47255\) octaves, amplitude \(0.011054\), and in-sample
\(R^2=0.643529\).  The complete causal forecast did not consistently improve
on a constant predictor.

The separate spectral coefficient estimates are
\begin{equation}
\begin{array}{c|cc}
K&\text{direct}&\text{cross-fitted}\\ \hline
14&0.11802&0.11793\\
50&0.10508&0.10508\\
100&0.10250&0.10270\\
102&0.10211&0.10231
\end{array}.
\label{eq:spectral-coefficient-table}
\end{equation}
No equality between these spectral coefficients and
\eqref{eq:density-normalized-quadratic-mass} is asserted.

\subsection{Excursion diagnostics}
\label{subsec:numerical-excursions}

At \(N=5\cdot10^6\),
\begin{equation}
m=3{,}039{,}633,\qquad
p=1{,}519{,}462,\qquad q=1{,}520{,}171.
\label{eq:five-million-sign-counts}
\end{equation}
There are 6,016 completed excursions in the full support-time path, so
\begin{equation}
\log_{10}|\cO_N(\mu)|
=6016\log_{10}2=1810.996\ldots.
\label{eq:reflection-orbit-decimal-size}
\end{equation}
A windowed analysis excluding excursions meeting \(n\leq1000\) contains
5,956 completed excursions.

A reported sample of 300 uniform bridges gave
\begin{equation}
\overline Z=1679.8,\qquad s_Z=1102.4,\qquad
\min Z=7,\qquad\max Z=5544,
\label{eq:return-sample-summary}
\end{equation}
so none reached the observed value 6,016.  The corresponding exact one-sided
95-percent zero-exceedance bound is
\begin{equation}
1-0.05^{1/300}=0.00994.
\label{eq:three-hundred-zero-bound}
\end{equation}
For the longest completed-excursion fraction, the same sample had mean
\(0.5301\), median \(0.4806\), and minimum \(0.1219\), compared with the
observed \(0.11737\).

\begin{table}[htbp]
\centering
\caption{Observed excursion-length tails and the unconditioned first-return
benchmark \(6016\sqrt{2/\pi}\,c^{-1/2}\).}
\label{tab:excursion-tail}
\begin{tabular}{r r r r}
\toprule
threshold \(c\)&observed&benchmark&ratio\\
\midrule
10&1474&1517.9&0.971\\
100&454&480.0&0.946\\
1000&140&151.8&0.922\\
10,000&40&48.0&0.833\\
\bottomrule
\end{tabular}
\end{table}

\begin{table}[htbp]
\centering
\caption{Large-cutoff support and return census.  The complete-cutoff return
count was not recorded in the cited census.}
\label{tab:large-cutoff-census}
\begin{tabular}{r r r r}
\toprule
\(N\)&nonzero entries&returns to zero&\(M(N)\)\\
\midrule
50,000,000&30,396,344&15,756&\(-908\)\\
100,000,000&60,792,694&25,525&1,928\\
150,000,000&91,189,087&34,232&1,089\\
200,000,000&121,585,426&36,069&320\\
250,000,000&151,981,736&36,755&\(-982\)\\
300,000,000&182,378,126&42,505&\(-592\)\\
350,000,000&212,774,472&47,215&2,896\\
455,052,511&276,638,778&not recorded&\(-5,414\)\\
\bottomrule
\end{tabular}
\end{table}
Large computations of the Mertens function provide numerical context for
these finite ranges \cite{DelegliseRivat1996,KotnikLune2004,Hurst2018}.

\subsection{Partition and conditional-choice diagnostics}
\label{subsec:numerical-partition-anova}

\begin{table}[htbp]
\centering
\caption{Partition audit.  Each \(p\)-value is two-sided; the Pearson
statistic has approximately one fewer degree of freedom than the number of
bins.}
\label{tab:partition-audit}
\begin{tabular}{r c c c c}
\toprule
bins&unweighted: observed/null&\(p\)&Pearson: observed/null&\(p\)\\
\midrule
12&0.252/0.352&0.135&6.07/11.06&0.188\\
24&0.351/0.512&0.0066&15.65/23.24&0.103\\
48&0.706/0.710&0.837&45.89/47.56&0.739\\
\bottomrule
\end{tabular}
\end{table}

\begin{table}[htbp]
\centering
\caption{Held-out audit of the local quadratic-only conditional-choice
model.  Each scale used 1,000 choice sets with 13 paths.}
\label{tab:local-anova}
\begin{tabular}{r c c c}
\toprule
\(q\)&deviance \(p\)&test gain per set&test top--1\\
\midrule
256&0.550&\(1.0\cdot10^{-6}\)&0.120\\
512&0.904&\(-4.16\cdot10^{-4}\)&0.040\\
1024&0.661&0.00245&0.120\\
\bottomrule
\end{tabular}
\end{table}

\subsection{Blind frequency audit}
\label{subsec:numerical-blind-frequencies}

\begin{table}[htbp]
\centering
\scriptsize
\caption{External audit of the first fourteen zero ordinates against the
pre-null blind spectral candidates.  Reference values were not used during
fitting.  A dash denotes no matched candidate.}
\label{tab:blind-frequency-audit}
\begin{tabular}{r r r r}
\toprule
mode&reference&blind estimate&error\\
\midrule
1&14.135&14.138&0.00347\\
2&21.022&21.017&\(-0.00461\)\\
3&25.011&25.011&0.00048\\
4&30.425&30.430&0.00471\\
5&32.935&32.931&\(-0.00422\)\\
6&37.586&37.559&\(-0.02737\)\\
7&40.919&--&--\\
8&43.327&--&--\\
9&48.005&48.015&0.00970\\
10&49.774&49.837&0.06298\\
11&52.970&--&--\\
12&56.446&--&--\\
13&59.347&59.377&0.03029\\
14&60.832&60.810&\(-0.02195\)\\
\bottomrule
\end{tabular}
\end{table}

\section{Reproducibility specification}
\label{app:reproducibility}

This appendix specifies the complete deposit required to reproduce every
computational statement in the paper.  It is a protocol, not a claim that an
independent rerun was performed while preparing this version.  The exact
finite identities can be checked without the data; the decimal tables require
the identified inputs and the declared software environment.

\subsection{Required archive structure}
\label{subsec:archive-structure}

The reproducibility archive shall contain the following top-level items:
\begin{enumerate}[label=(\roman*)]
\item \texttt{README.md}, with the single-command instructions, expected wall
time, memory requirement, and a map from each table and equation in the paper
to its raw output;
\item \texttt{MANIFEST.sha256}, containing a SHA-256 digest for every input,
source file, raw output, and generated table;
\item \texttt{environment.txt}, recording operating system, processor,
available memory, MATLAB release and update, BLAS library, thread count, and
locale;
\item \texttt{data/moebius.bin}, the signed-byte M\"obius word;
\item \texttt{data/zeta\_zeros\_400.csv} and
\texttt{data/zeta\_derivatives\_400.csv}, used only by the separate spectral
branch;
\item \texttt{src/}, containing the canonical source files listed below;
\item \texttt{raw/}, containing unmodified console logs, CSV files, and MAT
files from every run;
\item \texttt{tables/}, containing only deterministic transformations of the
raw files into manuscript tables;
\item \texttt{tests/}, containing the finite exhaustive tests and their
expected outputs.
\end{enumerate}
No source file may download data during a run.  The exact/permutation branch
and the spectral branch must have separate entry points, and the former must
abort if a spectral input is opened.

\subsection{Canonical source inventory}
\label{subsec:source-inventory}

The archive shall use mathematical names and contain, at minimum,
\begin{description}[leftmargin=4.3cm,style=nextline]
\item[\texttt{run\_all\_reproductions.m}]
Top-level orchestrator; runs tests, exact calculations, sampling experiments,
and table generation in a fixed order.
\item[\texttt{fixed\_mask\_moments.m}]
Streaming observed sum, exact class mean, exact variance, forward-error
audit, continuum diagnostics, and all fourteen cutoffs.
\item[\texttt{uniform\_bridge\_ladder.m}]
Uniform fixed-count sampler at the six Monte Carlo cutoffs.
\item[\texttt{excursion\_statistics.m}]
Completed-excursion decomposition, reflection count, exact return mean,
orientation randomization, Catalan controls, and length diagnostics.
\item[\texttt{partition\_and\_scale\_tests.m}]
Audits with 12, 24, and 48 bins.  Sliding scales, slow-envelope
correlations, split halves, and familywise calibration.
\item[\texttt{conditional\_choice\_analysis.m}]
Feature construction, conditional-logit models, chronological splits, and
deviance tables.
\item[\texttt{blind\_log\_spectrum.m}]
Blind discovery/validation/test spectral pipeline and the 100 complete
fixed-mask placebo reruns.
\item[\texttt{spectral\_second\_moment.m}]
Four-hundred-zero direct sum, tail alternatives, fourteen-cutoff
reconstruction, and residual decomposition.
\item[\texttt{make\_manuscript\_tables.m}]
Reads raw outputs and writes tables without recomputing any statistic.
\end{description}
Every function used by these entry points must be included as source rather
than supplied through an unrecorded local path.  The archive must also state
the license under which the sources and derived tables are released.

\subsection{Input definition and integrity checks}
\label{subsec:input-integrity}

The main data file has the immutable specification
\begin{equation}
\begin{split}
\text{format}&=\texttt{int8},\\
\text{allowed values}&=\{-1,0,1\},\\
\text{number of bytes}&=455{,}052{,}511,\\
\text{SHA-256}&=
\texttt{0e551f47d924ef6f3e348e131fdf53e0c677ad5e5471ef9d6249b3fd662e2ea4}.
\end{split}
\label{eq:main-input-specification}
\end{equation}
Before any statistic is evaluated, the program must:
\begin{enumerate}[label=(\alph*)]
\item verify the byte count, allowed alphabet, and SHA-256 digest;
\item independently sieve \(\mu(n)\) for \(1\leq n\leq10^6\) and compare
every entry;
\item verify at each cutoff the nonzero count and endpoint recorded in the
raw reference table;
\item abort unless
\begin{equation}
M(N)\equiv Q(N)\pmod2
\label{eq:reproducibility-parity-test}
\end{equation}
at every cutoff;
\item verify \(p=(m+S)/2\), \(q=(m-S)/2\), and \(T_m=S\).
\end{enumerate}

The spectral CSV files must contain one row per positive ordinate with columns
\texttt{index}, \texttt{gamma}, \texttt{ReZetaPrime},
\texttt{ImZetaPrime}, \texttt{AbsZetaPrime}, working precision, software
used, and an independent residual or interval check for
\(\zeta(1/2+i\gamma)\).  Their hashes must be listed in
\texttt{MANIFEST.sha256}.  The code must recompute
\(|\zeta'(\rho)|\) from its real and imaginary parts and abort if it disagrees
with the stored absolute value beyond the declared tolerance.

\subsection{Exact-layer configuration}
\label{subsec:exact-layer-configuration}

The main run uses
\begin{align}
N_0&=1000,
\label{eq:reproducibility-n-zero}\\
\mathcal N_{\mathrm{exact}}
&=\{10^4,3\cdot10^4,10^5,3\cdot10^5,10^6,5\cdot10^6,
\notag\\
&\hspace{8mm}5\cdot10^7,10^8,1.5\cdot10^8,2\cdot10^8,
2.5\cdot10^8,3\cdot10^8,3.5\cdot10^8,455052511\}.
\label{eq:reproducibility-exact-cutoffs}
\end{align}
The streaming read block is \(2^{22}\) bytes and the positive-sum audit
block is \(2^{16}\) terms.  Products that can exceed
\(2^{53}\) must not be formed.  In particular, the implementation shall use
\begin{equation}
\left(\frac{M(n)}n\right)^2,\qquad
\frac{Q(n)}n\frac{Q(n)-1}{n},\qquad
\frac{Q(n)}n\frac1n
\label{eq:safe-division-forms}
\end{equation}
rather than forming \(M(n)^2\), \(Q(n)^2\), or \(n^2\) first at the largest
cutoffs.

The mean must be computed independently in two ways:
\begin{enumerate}[label=(\alph*)]
\item from the streaming sums \(\Sigma_1,\Sigma_2\) in
\eqref{eq:sigma-one}--\eqref{eq:sigma-two};
\item from the explicit plateau weights and
\eqref{eq:exact-class-mean-alternative}.
\end{enumerate}
The variance must likewise be evaluated by the prefix formula
\eqref{eq:exact-class-variance} and, at \(N=10^4\), by the literal triangular
\(O(m^2)\) sum \eqref{eq:variance-triangular-sum}.

\subsection{Floating-point audit}
\label{subsec:floating-point-audit}

The recorded main environment is MATLAB
24.2.0.2740171 (R2024b Update 1).  The forward-error argument assumes
IEEE-754 binary64, round-to-nearest ties-to-even, no extended precision, no
fused-operation contraction, and no reassociation.  If the execution
environment does not satisfy these conditions, the program must label the
audit inapplicable or replace it with directed-rounding interval arithmetic.

Within each block of \(r\) positive summands, ordinary summation is bounded
using
\begin{equation}
\gamma_r=\frac{ru}{1-ru},\qquad u=2^{-53},
\label{eq:higham-gamma}
\end{equation}
and block totals are accumulated with compensated summation.  Bounds are
propagated through the algebra for \(\overline{\QF}_N\) and
\(\Delta_N=\QF(\mu)-\overline{\QF}_N\).  The run must abort unless
\begin{equation}
0\notin[\Delta_N^{-},\Delta_N^{+}]
\quad\text{and}\quad \Delta_N^{+}<0
\label{eq:separation-abort-condition}
\end{equation}
at every cutoff.  This procedure is an a-priori forward-error audit in the
sense of \cite{Higham2002}; it is not interval arithmetic with directed
rounding.

\subsection{Mandatory exact self-tests}
\label{subsec:mandatory-self-tests}

The test suite must run before the data file is processed and must abort on
any failure.  It includes:
\begin{enumerate}[label=(\roman*)]
\item complete enumeration of every fixed-count sign class with
\(1\leq m\leq8\), covering 56 classes and 1,890 covariance pairs;
\item comparison of every enumerated value
\(\Cov(T_j^2,T_k^2)\) with
\eqref{eq:factored-covariance}; the recorded maximum absolute error is
\(3.2\cdot10^{-13}\);
\item explicit enumeration of the 20 words with \(m=6,p=q=3\), including
\(\mu_2=-1/5\), \(\mu_4=1/5\), and
\begin{equation}
\Var(T_2^2)=3.84;
\label{eq:m-six-variance-test}
\end{equation}
\item checks that \(\psi(m)=0\) and
\(\E T_m^2=S^2\) identically;
\item equality of the \(O(m)\) and \(O(m^2)\) variance implementations at
\(N=10^4\), with recorded relative discrepancy \(3.05\cdot10^{-16}\);
\item equality of the streaming and plateau-weight mean and variance routes,
with recorded mismatches between \(7\cdot10^{-18}\) and \(10^{-16}\);
\item exhaustive checks of Proposition \ref{prop:exact-return-mean} for at
least \((p,q)=(3,3),(4,3),(5,4),(6,4)\);
\item a SHA-256 unit test against a published standard test vector;
\item a symbolic or high-precision numerical self-test of the Legendre
transform in \eqref{eq:ou-rate-function}, explicitly labelled as a transform
test rather than a proof of the OU transfer.
\end{enumerate}

The reference arithmetic printed by the test suite is
\begin{align}
1-0.05^{1/2800}&=0.00106933223978,
\label{eq:reference-upper-2800}\\
(1-0.00428)^{2800}&=6.08478495942\cdot10^{-6},
\label{eq:reference-zero-2800}\\
1-0.05^{1/9500}&=0.000315290524814,
\label{eq:reference-upper-9500}\\
9500(0.00428)&=40.66.
\label{eq:reference-expected-9500}
\end{align}

\subsection{Uniform sampling protocol}
\label{subsec:uniform-sampling-protocol}

The Monte Carlo configuration is
\begin{equation}
\begin{array}{c|rrrrrr}
N&10^4&3\cdot10^4&10^5&3\cdot10^5&10^6&5\cdot10^6\\ \hline
R&300000&150000&100000&50000&10000&10000\\
\text{seed}&19092021&19092022&19092023&19092024&19092025&19092026
\end{array}.
\label{eq:monte-carlo-configuration}
\end{equation}
The generator is MT19937, invoked as
\texttt{rng(seed,'twister')} separately at each cutoff.  A draw is produced
by a uniform random permutation of the \(p\) positive and \(q\) negative
signs, followed by placement on the fixed support.  The program must check
\(T_m=S\) on every draw and test the marginal positive-sign frequency at
fixed support indices.

For a positive hit count, the archive reports the Wilson interval and the
exact binomial likelihood.  For zero hits, it reports the one-sided exact
bound
\begin{equation}
p_{\mathrm{upper}}=1-\alpha^{1/R},\qquad \alpha=0.05.
\label{eq:clopper-pearson-zero-bound}
\end{equation}
The sample standard-deviation check uses
\begin{equation}
\operatorname{SE}(s)
=\frac{s}{2}\sqrt{\frac{\widehat\kappa-1}{R}},
\qquad
\widehat\kappa=\frac{\widehat m_4}{s^4}.
\label{eq:sample-standard-deviation-se}
\end{equation}
The program must print
\((s-\sigma_N)/\operatorname{SE}(s)\) and require its absolute value to be
below the declared calibration threshold.

The two-moment comparison uses
\begin{equation}
\nu=\frac{2\overline{\QF}_N^2}{\sigma_N^2},\qquad
\QF_{\chi^2}=\frac{\overline{\QF}_N}{\nu}\chi_{\nu}^2.
\label{eq:two-moment-chi-square-model}
\end{equation}
It is a deliberately falsifiable approximation, not the sampling law.
Every sampled quadratic value, not only summary statistics, must be written
to a per-cutoff CSV file.

MT19937 has 19,937 state bits and cannot visit every member of the largest
class.  The archive shall state this limitation.  The exact mean and variance
are unaffected; the Monte Carlo results concern the implemented uniform
finite-permutation algorithm and require an independent generator as a
robustness check before publication.

\subsection{Excursion and conditional-choice protocols}
\label{subsec:secondary-protocols}

The excursion archive must include:
\begin{enumerate}[label=(\alph*)]
\item the complete ordered list of return times and completed-excursion
lengths;
\item a flag showing whether the terminal segment is incomplete and excluded
from reflection;
\item 100,000 orientation-null realizations with the unsigned skeleton fixed;
\item the 500 observed excursions and all 100 Catalan replicates used for the
shape tests;
\item all 300 uniform bridges used for the return-count and longest-excursion
summaries;
\item the exact seed and generator state for each component.
\end{enumerate}

The conditional-choice archive must include all 1,000 choice sets with 13
paths at each \(q=256,512,1024\), the feature matrix before and after
standardization, the chronological split indices, coefficient estimates,
standard errors, log likelihoods, deviance increments, degrees of freedom,
held-out log scores, and false-discovery adjustment.  Expected cell counts
must be reported for every chi-square diagnostic; where they are too small,
an exact or permutation calibration is required.

\subsection{Blind spectral protocol}
\label{subsec:blind-spectral-protocol}

The blind log-spectrum run uses \(N=5\cdot10^6\), 1,321 logarithmic samples,
frequency ceiling \(216.275\), four hierarchical levels, and growing capacity
\begin{equation}
J(n)=\lfloor n^{0.38}\rfloor,\qquad J_{\max}=15.
\label{eq:blind-spectrum-capacity}
\end{equation}
The chronological discovery, validation, and test endpoints are
\(0.45N,0.65N,N\).  Reference ordinates must not be loaded until the blind
fit and all 100 placebo fits are complete.  Each placebo must rerun the
entire selection pipeline, including exponent choice, peak selection,
frequency refinement, and validation.  The archive must preserve:
\begin{enumerate}[label=(\alph*)]
\item all 1,321 sampled log-time values and observed signals;
\item every candidate peak before and after hierarchical selection;
\item all 100 placebo candidate sets and test gains;
\item the mode-by-mode multiple-testing table;
\item the external reference audit generated only after selection;
\item warnings, ranks, and condition numbers from every linear solve.
\end{enumerate}
The pre-null match table is an algorithmic recovery result.  The recorded
100-placebo run is a pilot resolution audit only: by
\eqref{eq:blind-bh-resolution-obstruction}, zero surviving modes cannot be
used as confirmatory evidence of absence.  A confirmatory modewise analysis
at \(q=0.010\) must rerun the complete pipeline on at least 1,099 placebos,
and should use more if stable tail estimation is required.  Alternatively,
the protocol may prespecify one global maximum statistic and calibrate that
single statistic with the full-pipeline placebo distribution.  The choice
between these two designs must be made before a new run and recorded in the
manifest.

\subsection{Separate spectral-second-moment protocol}
\label{subsec:spectral-second-moment-protocol}

The spectral run must use arbitrary precision high enough that rounding in
each term of \eqref{eq:four-hundred-zero-sum} is negligible at the printed
precision.  It must output:
\begin{enumerate}[label=(\alph*)]
\item all 400 individual values
\(2/|\rho_k\zeta'(\rho_k)|^2\) and their cumulative sums;
\item the contributions of the first, first five, and first one hundred
zeros;
\item the precise last-quarter estimator defining \(\widehat c=0.0923\);
\item both tail alternatives \(\widehat c\) and \(3/\pi^3\);
\item the fourteen reconstructed quadratic-mass values, their residuals, and
the zero-parameter RMS;
\item the fit \(R_N=A+B/L_N\), its covariance matrix, and the discrete
sum--integral correction;
\item all working-precision and zero-validation diagnostics.
\end{enumerate}
No output from this branch may be merged with the no-zero exact run before
both raw manifests have been written.

\subsection{Raw output inventory and acceptance criteria}
\label{subsec:raw-output-inventory}

At minimum, the archive shall preserve
\begin{equation}
\begin{gathered}
\texttt{exact\_layer.csv},\quad
\texttt{permutation\_ladder.csv},\quad
\texttt{sampled\_quadratic\_values.csv},\\
\texttt{continuum\_diagnostics.csv},\quad
\texttt{ou\_quantiles.csv},\quad
\texttt{excursions.csv},\\
\texttt{partition\_audit.csv},\quad
\texttt{conditional\_choice.csv},\quad
\texttt{blind\_modes.csv},\\
\texttt{blind\_placebos.csv},\quad
\texttt{spectral\_terms.csv}.
\end{gathered}
\label{eq:required-output-files}
\end{equation}
The MAT files must retain every accumulator and configuration object needed
to audit the CSV files.  Console output must be captured verbatim.

A reproduction is accepted only if:
\begin{enumerate}[label=(\roman*)]
\item every input and source hash matches the manifest;
\item all exhaustive and integrity tests pass;
\item the fourteen observed values, exact means, and exact variances reproduce
within the stated forward-error bounds;
\item the sign of every separation interval is negative;
\item raw Monte Carlo summaries reproduce from the declared seeds, and an
independent generator gives compatible confidence intervals;
\item the numerical tables are generated only from raw files;
\item every table cell in the manuscript is traceable to one raw column and
row;
\item a fresh LaTeX build completes without undefined references, missing
citations, or overfull material that obscures data.
\end{enumerate}

The intended single command is
\begin{equation}
\texttt{matlab -batch "run\_all\_reproductions"}
\label{eq:single-reproduction-command}
\end{equation}
from the archive root.  The recorded full exact and Monte Carlo run required
2,133.7 seconds in the stated environment.  The command and all programs in
this appendix are specified for an independent rerun; they were not executed
as part of the editorial production of Version 3.

\begin{thebibliography}{99}
\raggedright

\bibitem{BardenetMaillard2015}
R. Bardenet and O.-A. Maillard,
\newblock Concentration inequalities for sampling without replacement,
\newblock \emph{Bernoulli} \textbf{21} (2015), 1361--1385,
\newblock \href{https://doi.org/10.3150/14-BEJ605}
{doi:10.3150/14-BEJ605}, arXiv:1309.4029.

\bibitem{BenjaminiHochberg1995}
Y. Benjamini and Y. Hochberg,
\newblock Controlling the false discovery rate: a practical and powerful
approach to multiple testing,
\newblock \emph{Journal of the Royal Statistical Society B} \textbf{57}
(1995), 289--300,
\newblock \href{https://doi.org/10.1111/j.2517-6161.1995.tb02031.x}
{doi:10.1111/j.2517-6161.1995.tb02031.x}.

\bibitem{ClopperPearson1934}
C. J. Clopper and E. S. Pearson,
\newblock The use of confidence or fiducial limits illustrated in the case
of the binomial,
\newblock \emph{Biometrika} \textbf{26} (1934), 404--413,
\newblock \href{https://doi.org/10.1093/biomet/26.4.404}
{doi:10.1093/biomet/26.4.404}.

\bibitem{Cochran1977}
W. G. Cochran,
\newblock \emph{Sampling Techniques}, 3rd ed.,
\newblock Wiley, New York, 1977.

\bibitem{DelegliseRivat1996}
M. Del\'eglise and J. Rivat,
\newblock Computing the summation of the M\"obius function,
\newblock \emph{Experimental Mathematics} \textbf{5} (1996), 291--295,
\newblock \href{https://doi.org/10.1080/10586458.1996.10504365}
{doi:10.1080/10586458.1996.10504365}.

\bibitem{DonskerVaradhan1975}
M. D. Donsker and S. R. S. Varadhan,
\newblock Asymptotic evaluation of certain Markov process expectations for
large time, I,
\newblock \emph{Communications on Pure and Applied Mathematics} \textbf{28}
(1975), 1--47,
\newblock \href{https://doi.org/10.1002/cpa.3160280102}
{doi:10.1002/cpa.3160280102}.

\bibitem{Feller1968}
W. Feller,
\newblock \emph{An Introduction to Probability Theory and Its Applications},
Vol.~I, 3rd ed.,
\newblock Wiley, New York, 1968.

\bibitem{FlajoletSedgewick2009}
P. Flajolet and R. Sedgewick,
\newblock \emph{Analytic Combinatorics},
\newblock Cambridge University Press, Cambridge, 2009,
\newblock \href{https://doi.org/10.1017/CBO9780511801655}
{doi:10.1017/CBO9780511801655}.

\bibitem{Gonek1989}
S. M. Gonek,
\newblock On negative moments of the Riemann zeta-function,
\newblock \emph{Mathematika} \textbf{36} (1989), 71--88,
\newblock \href{https://doi.org/10.1112/S0025579300013589}
{doi:10.1112/S0025579300013589}.

\bibitem{HartmanWintner1941}
P. Hartman and A. Wintner,
\newblock On the law of the iterated logarithm,
\newblock \emph{American Journal of Mathematics} \textbf{63} (1941),
169--176,
\newblock \href{https://doi.org/10.2307/2371287}
{doi:10.2307/2371287}.

\bibitem{Higham2002}
N. J. Higham,
\newblock \emph{Accuracy and Stability of Numerical Algorithms}, 2nd ed.,
\newblock SIAM, Philadelphia, 2002,
\newblock \href{https://doi.org/10.1137/1.9780898718027}
{doi:10.1137/1.9780898718027}.

\bibitem{Hoeffding1951}
W. Hoeffding,
\newblock A combinatorial central limit theorem,
\newblock \emph{The Annals of Mathematical Statistics} \textbf{22} (1951),
558--566,
\newblock \href{https://doi.org/10.1214/aoms/1177729545}
{doi:10.1214/aoms/1177729545}.

\bibitem{Hoeffding1963}
W. Hoeffding,
\newblock Probability inequalities for sums of bounded random variables,
\newblock \emph{Journal of the American Statistical Association}
\textbf{58} (1963), 13--30,
\newblock \href{https://doi.org/10.1080/01621459.1963.10500830}
{doi:10.1080/01621459.1963.10500830}.

\bibitem{HubertSchultz1976}
L. Hubert and J. Schultz,
\newblock Quadratic assignment as a general data analysis strategy,
\newblock \emph{British Journal of Mathematical and Statistical Psychology}
\textbf{29} (1976), 190--241,
\newblock \href{https://doi.org/10.1111/j.2044-8317.1976.tb00714.x}
{doi:10.1111/j.2044-8317.1976.tb00714.x}.

\bibitem{Hurst2018}
G. Hurst,
\newblock Computations of the Mertens function and improved bounds on the
Mertens conjecture,
\newblock \emph{Mathematics of Computation} \textbf{87} (2018), 1013--1028,
\newblock \href{https://doi.org/10.1090/mcom/3277}
{doi:10.1090/mcom/3277}, arXiv:1610.08551.

\bibitem{KotnikLune2004}
T. Kotnik and J. van de Lune,
\newblock Further systematic computations on the summatory function of the
M\"obius function,
\newblock \emph{Experimental Mathematics} \textbf{13} (2004), 449--453,
\newblock \href{https://doi.org/10.1080/10586458.2004.10504549}
{doi:10.1080/10586458.2004.10504549}.

\bibitem{McFadden1974}
D. McFadden,
\newblock Conditional logit analysis of qualitative choice behavior,
\newblock in P. Zarembka (ed.), \emph{Frontiers in Econometrics},
Academic Press, New York, 1974, pp.~105--142.

\bibitem{MontgomeryVaughan2007}
H. L. Montgomery and R. C. Vaughan,
\newblock \emph{Multiplicative Number Theory I: Classical Theory},
\newblock Cambridge University Press, Cambridge, 2007,
\newblock \href{https://doi.org/10.1017/CBO9780511618314}
{doi:10.1017/CBO9780511618314}.

\bibitem{Ng2004}
N. Ng,
\newblock The distribution of the summatory function of the M\"obius
function,
\newblock \emph{Proceedings of the London Mathematical Society}
\textbf{89} (2004), 361--389,
\newblock \href{https://doi.org/10.1112/S0024611504014741}
{doi:10.1112/S0024611504014741}, arXiv:math/0310381.

\bibitem{OdlyzkoTeRiele1985}
A. M. Odlyzko and H. J. J. te Riele,
\newblock Disproof of the Mertens conjecture,
\newblock \emph{Journal f\"ur die reine und angewandte Mathematik}
\textbf{357} (1985), 138--160,
\newblock \href{https://doi.org/10.1515/crll.1985.357.138}
{doi:10.1515/crll.1985.357.138}.

\bibitem{Sen1972}
P. K. Sen,
\newblock Finite population sampling and weak convergence to a Brownian
bridge,
\newblock \emph{Sankhy\=a, Series A} \textbf{34} (1972), 85--90,
\newblock MR0345270, JSTOR 25049775.

\bibitem{Serfling1974}
R. J. Serfling,
\newblock Probability inequalities for the sum in sampling without
replacement,
\newblock \emph{The Annals of Statistics} \textbf{2} (1974), 39--48,
\newblock \href{https://doi.org/10.1214/aos/1176342611}
{doi:10.1214/aos/1176342611}.

\bibitem{Stanley1999}
R. P. Stanley,
\newblock \emph{Enumerative Combinatorics}, Vol.~2,
\newblock Cambridge University Press, Cambridge, 1999,
\newblock \href{https://doi.org/10.1017/CBO9780511609589}
{doi:10.1017/CBO9780511609589}.

\bibitem{Titchmarsh1986}
E. C. Titchmarsh,
\newblock \emph{The Theory of the Riemann Zeta-Function}, 2nd ed.,
\newblock revised by D. R. Heath-Brown,
\newblock Oxford University Press, Oxford, 1986.

\bibitem{Wilson1927}
E. B. Wilson,
\newblock Probable inference, the law of succession, and statistical
inference,
\newblock \emph{Journal of the American Statistical Association}
\textbf{22} (1927), 209--212,
\newblock \href{https://doi.org/10.1080/01621459.1927.10502953}
{doi:10.1080/01621459.1927.10502953}.

\bibitem{ZhouEtAl2013}
Y.-H. Zhou, G. Mayhew, Z. Sun, X. Xu, F. Zou, and F. A. Wright,
\newblock Space--time clustering and the permutation moments of quadratic
forms,
\newblock \emph{Stat} \textbf{2} (2013), 292--302,
\newblock \href{https://doi.org/10.1002/sta4.37}
{doi:10.1002/sta4.37}.

\end{thebibliography}

\end{document}
